\documentclass[10pt]{article}
\usepackage[utf8]{inputenc}
\usepackage[T1]{fontenc}
\usepackage{amsmath}
\usepackage{amsfonts}
\usepackage{amssymb}
\usepackage[version=4]{mhchem}
\usepackage{stmaryrd}
\usepackage{bbold}
\usepackage{hyperref}
\hypersetup{colorlinks=true, linkcolor=blue, filecolor=magenta, urlcolor=cyan,}
\urlstyle{same}
\usepackage{caption}
\usepackage{multirow}

\title{Biprojective Almost Perfect Nonlinear Functions }

\author{Faruk Göloğlu ${ }^{\text {® }}$}
\date{}


%New command to display footnote whose markers will always be hidden
\let\svthefootnote\thefootnote
\newcommand\blfootnotetext[1]{%
  \let\thefootnote\relax\footnote{#1}%
  \addtocounter{footnote}{-1}%
  \let\thefootnote\svthefootnote%
}

%Overriding the \footnotetext command to hide the marker if its value is `0`
\let\svfootnotetext\footnotetext
\renewcommand\footnotetext[2][?]{%
  \if\relax#1\relax%
    \ifnum\value{footnote}=0\blfootnotetext{#2}\else\svfootnotetext{#2}\fi%
  \else%
    \if?#1\ifnum\value{footnote}=0\blfootnotetext{#2}\else\svfootnotetext{#2}\fi%
    \else\svfootnotetext[#1]{#2}\fi%
  \fi
}

\DeclareUnicodeCharacter{2020}{\ifmmode\dagger\else{$\dagger$}\fi}

\begin{document}
\maketitle
\captionsetup{singlelinecheck=false}


\begin{abstract}
In this paper, we introduce the concept of biprojectivity in studying the cryptographically important almost perfect nonlinear (APN) functions. Although several known families of biprojective APN functions exist in the literature, the concept itself has not been explicitly observed and studied in detail. We give a survey of known APN families and functions that fall into the biprojective setting, then provide a method for finding such families and functions. We finally give two new infinite families of biprojective APN functions CCZ-inequivalent to any known APN function and study their properties.
\end{abstract}

Index Terms-APN functions, projective polynomials.

\section*{I. Introduction}
IN THIS paper, we introduce biprojective almost perfect nonlinear (APN) functions. After introducing the concept, we introduce a method for finding such functions. We then proceed to give two infinite families of biprojective APN functions that are CCZ-inequivalent to any known APN function (Theorem III.2). We also show that many known infinite families of APN functions fall into the biprojective setting. Importantly, the Kim function [1], which is the only known APN function equivalent to an APN permutation on an even dimensional $\mathbb{F}_{2}$-vector space, falls into the biprojective setting. Finding APN permutations on even dimensions is regarded as one of the most important problems in this research area. Thus, finding infinite families of biprojective APN functions is important in this regard, even though there exist several (more precisely, thirteen, excluding the ones found in this paper) known infinite families of quadratic APN functions. We also find a sporadic quadratic APN function that satisfies a property (related to being equivalent to a permutation) which is shared only by the Gold infinite family and the sporadic Kim function among the known quadratic APN functions (Proposition III.4).

In Section II, we explain some basic notions. In Section III, we define biprojective functions and give a brief survey on known biprojective APN families. In Section IV, we introduce our method. In Section V, we prove the APN property of the new families. Finally, in Section VI, we prove that the introduced APN families are indeed new.

\section*{II. Preliminaries}
A function $F: \mathbb{F}_{2}^{n} \rightarrow \mathbb{F}_{2}^{l}$ is called a vectorial Boolean function (or an $(n, l)$-function). When $l=1$ then we simply

\footnotetext{Manuscript received June 16, 2020; revised February 15, 2022; accepted February 20, 2022. Date of publication March 8, 2022; date of current version June 15, 2022. This work was supported by the GAČR Grant 18-19087S-301-13/201843.\\
The author is with the Department of Mathematics, Faculty of Mathematics and Physics, Charles University, Praha, 18675 Prague, Czech Republic (e-mail: \href{mailto:faruk.gologlu@mff.cuni.cz}{faruk.gologlu@mff.cuni.cz}).

Communicated by S. Mesnager, Associate Editor for Sequences.\\
Digital Object Identifier 10.1109/TIT.2022.3157798
}
say that $F$ is a Boolean function. In this paper, we only deal with the $l=n$ case. The finite field $\mathbb{F}_{2^{n}}$ is the degree $n$ extension of the finite field $\mathbb{F}_{2}=\{0,1\}$. The multiplicative group of $\mathbb{F}_{2^{n}}$ is denoted by $\mathbb{F}_{2^{n}}^{\times}=\mathbb{F}_{2^{n}} \backslash\{0\}$. We will identify the vector space $\mathbb{F}_{2}^{n}$ with the additive structure of the finite field $\mathbb{F}_{2^{n}}$. Thus $F$ can be viewed as a mapping from $\mathbb{F}_{2^{n}}$ to itself. Using Lagrange interpolation, any vectorial Boolean function $F$ can be written uniquely as a polynomial in the polynomial ring $\mathbb{F}_{2^{n}}[X]$ with (polynomial) degree at most $2^{n}-1$, i.e.,
$$
F(X)=\sum_{0 \leq i<2^{n}} A_{i} X^{i}, \quad A_{i} \in \mathbb{F}_{2^{n}}
$$
which is called the univariate representation of $F$. We will not make a distinction between the functions and their polynomial representations. Note that the (polynomial) degree is defined as $\max \left\{i: A_{i} \neq 0\right\}$. We define the algebraic degree of a vectorial Boolean function $F$ as $\max \left\{\mathrm{wt}_{2}(i): A_{i} \neq 0\right\}$ where $\mathrm{wt}_{2}(i)$ denotes the number of ones in the base2 representation of the integer $i$. The linear and quadratic functions are defined using the notion of algebraic degree. Recall that a map $L: \mathbb{F}_{2^{n}} \rightarrow \mathbb{F}_{2^{n}}$ is said to be linear if $L$ is an $\mathbb{F}_{2}$-linear vector space endomorphism of $\mathbb{F}_{2^{n}}$. The univariate representation of a linear map $L$ is of the form
$$
L(X)=\sum_{0 \leq i<n} B_{i} X^{2^{i}}, \quad B_{i} \in \mathbb{F}_{2^{n}}
$$
which is called a linearized polynomial and its algebraic degree is at most one. An affine map $A: \mathbb{F}_{2^{n}} \rightarrow \mathbb{F}_{2^{n}}$ is of the form $A(X)=L(X)+C$ with $L$ linear and $C \in \mathbb{F}_{2^{n}}$. The quadratic functions are the polynomials
$$
Q(X)=\sum_{0 \leq i \neq j<n} A_{i j} X^{2^{i}+2^{j}}+L(X)+C, \quad A_{i j} \in \mathbb{F}_{2^{n}},
$$
where $L$ is linear and $C \in \mathbb{F}_{2^{n}}$ is a constant. A vectorial Boolean function $F: \mathbb{F}_{2^{n}} \rightarrow \mathbb{F}_{2^{n}}$ is called almost perfect nonlinear (APN), if for all $A, B \in \mathbb{F}_{2}^{n}$, with $A \neq 0$, the equation
$$
F(X)+F(X+A)=B
$$
has zero or two solutions. If $F$ is quadratic, then the APN condition is simplified as follows: $F$ is APN if and only if

\begin{equation*}
F(X+A)+F(X)+F(A)+F(0)=0 \tag{1}
\end{equation*}

for every $A \in \mathbb{F}_{2^{n}}^{\times}$, has exactly two solutions. Two vectorial Boolean functions $F$ and $G$ are said to be extended affine equivalent (EA-equivalent) if
$$
G(X)=A_{1} \circ F \circ A_{2}(X)+A_{3}(X),
$$
for affine permutations $A_{1}, A_{2}$ and an affine map $A_{3}$, and we write $F \sim_{\text {eA }} G$. We define the graph of $F$ as
$$
\Gamma_{F}=\left\{(X, F(X)) \quad X \in \mathbb{F}_{2^{n}}\right\} .
$$

Then, $F$ and $G$ are said to be graph equivalent (or CCZ-equivalent) [2] if

$$
\mathcal{A}\left(\Gamma_{F}\right)=\Gamma_{G},
$$

where $\mathcal{A}$ is an affine permutation of $\mathbb{F}_{2^{n}} \times \mathbb{F}_{2^{n}}$, and we write $F \sim_{\text {ccz }} G$. It has been shown that EA-equivalence and CCZ-equivalence coincide in the case of two quadratic APN functions [3]. A fundamental $\mathbb{F}_{2^{d}}$-linear map from an extension $\mathbb{F}_{2^{k d}}$ to $\mathbb{F}_{2^{d}}$ is the trace map:

$$
\operatorname{tr}_{\mathbb{F}_{2^{k d}} / \mathbb{F}_{2^{d}}}(X)=\sum_{i=0}^{k-1} X^{2^{d i}}
$$

For $X \in \mathbb{F}_{2^{n}}$, we will use

$$
\operatorname{tr}(X)=\operatorname{tr}_{\mathbb{F}_{2^{n}} / \mathbb{F}_{2}}(X),
$$

for the absolute trace. The trace- 0 hyperplane is defined to be the $(n-1)$-dimensional vector space

$$
H=\left\{x \in \mathbb{F}_{2^{n}} \operatorname{tr}(x)=0\right\} .
$$

The Walsh transform of $F: \mathbb{F}_{2^{n}} \rightarrow \mathbb{F}_{2^{n}}$ is defined as

$$
\widehat{F}(U, V)=\sum_{X \in \mathbb{F}_{2^{n}}}(-1)^{\operatorname{tr}(U F(X)+V X)}
$$

for $U, V \in \mathbb{F}_{2^{n}}$. Let $\mathrm{NB}_{F}$ be defined as

$$
\mathrm{NB}_{F}=\left\{U \in \mathbb{F}_{2^{n}} \widehat{F}(U, V)=0 \text { for some } V \in \mathbb{F}_{2^{n}}\right\} .
$$

Bent functions are the Boolean functions $f: \mathbb{F}_{2^{n}} \rightarrow \mathbb{F}_{2}$ whose Walsh transform satisfy $\widehat{f}(V)= \pm 2^{n / 2}$, where the Walsh transform (for single output Boolean functions) is defined as

$$
\widehat{f}(V)=\sum_{X \in \mathbb{F}_{2^{n}}}(-1)^{f(X)+\operatorname{tr}(V X)},
$$

for $V \in \mathbb{F}_{2^{n}}$. The Walsh transform of a quadratic Boolean function $f: \mathbb{F}_{2^{n}} \rightarrow \mathbb{F}_{2}$ satisfies $\widehat{f}(V) \in\left\{0, \pm 2^{s_{f}}\right\}$, where $s_{f}$ is an integer determined by $f$. A (not necessarily quadratic) Boolean function with this property is called a plateaued function. The component functions of $F$ are defined to be the Boolean functions $F_{U}(X)=\operatorname{tr}(U F(X))$ for $U \in \mathbb{F}_{2^{n}}$ (some authors prefer to exclude the zero function). If $F$ is a vectorial Boolean function all of whose component functions are plateaued, it is said to be a component-wise plateaued function. For a component-wise plateaued function $F$, the set $\mathrm{NB}_{F}$ coincides with the set of elements $U \in \mathbb{F}_{2^{n}}$ for which $F_{U}$ is not bent as a consequence of Parseval's identity which states that $\sum_{V \in \mathbb{F}_{2} n}(\widehat{f}(V))^{2}=2^{2 n}$. Since bent functions exist only when $n=2 m$, if $n$ is odd then $\mathrm{NB}_{F}=\mathbb{F}_{2^{n}}$. We define the vector

$$
\mathrm{N}_{F}=\left[\eta_{d}\left(\mathrm{NB}_{F}\right) 1 \leq d \leq n\right],
$$

where $\eta_{d}(S)$ is the number of $\mathbb{F}_{2}$-vector spaces of dimension $d$ in $S$. The sequence $\mathrm{N}_{F}$ has been shown to be an EAinvariant [4]. We will employ this invariant to distinguish APN families and exploit its relevance to the problem of finding APN permutations on even degree extensions. A vectorial Boolean function $F: \mathbb{F}_{2^{n}} \rightarrow \mathbb{F}_{2^{n}}$ is said to be a permutation if the map $x \mapsto F(x)$ permutes $\mathbb{F}_{2}{ }^{n}$. When $n$ is odd, there are\\
infinite families of APN permutations of $\mathbb{F}_{2}^{n}$. The monomials

$$
\mathrm{G}_{i}(X)=X^{2^{i}+1}
$$

are APN, whenever $\operatorname{gcd}(i, n)=1$ and are called Gold APN functions [5]. Note that the monomial mapping $X \mapsto X^{d}$ permutes $\mathbb{F}_{2^{n}}$ if and only if $\operatorname{gcd}\left(d, 2^{n}-1\right)=1$. Therefore, Gold APN functions are permutations if and only if $n$ is odd, since if $\operatorname{gcd}(i, n)=1$, then $\operatorname{gcd}\left(2^{i}+1,2^{n}-1\right)=1$ if and only if $n$ is odd. We also have $\operatorname{gcd}\left(2^{i}-1,2^{n}-1\right)= 2^{\operatorname{gcd}(i, n)}-1$. A major research problem is the existence of APN permutations of $\mathbb{F}_{2}^{n}$ when $n$ is even and larger than six. When $n \in\{2,4\}$, it has been shown by computer search that no APN permutation exists [6]. The only known example of an APN permutation on $\mathbb{F}_{2^{n}}$ with even $n$ is the case found by Browning et al. [1] when $n=6$. The permutation is CCZ-equivalent to the quadratic function

$$
\kappa(X)=X^{3}+X^{10}+U X^{24}
$$

on $\mathbb{F}_{2^{6}}[X]$, where $U$ is a root of $X^{6}+X^{4}+X^{3}+X+1$. See [7, Tables I, II] for the most recent list of known APN functions. Note that plateaued (and in particular quadratic) APN functions cannot be permutations of $\mathbb{F}_{2^{n}}$ when $n$ is even [8], [9] but can be CCZ-equivalent to a permutation as the above example shows. It is an interesting problem to exhibit other such cases.

We define the projective line $\mathbb{P}^{1}\left(\mathbb{F}_{2^{n}}\right)$ over the finite field $\mathbb{F}_{2^{n}}$ as the set of ratios $v / w$ where $(v, w) \in \mathbb{F}_{2^{n}} \times \mathbb{F}_{2^{n}} \backslash \{(0,0)\}$, with $v / 0$ for $v \in \mathbb{F}_{2^{n}}^{\times}$is defined to be the symbol $\infty$. Thus $\mathbb{P}^{1}\left(\mathbb{F}_{2^{n}}\right)=\mathbb{F}_{2^{n}} \cup\{\infty\}$. It is straightforward to show, when $A D+B C \neq 0$ for $A, B, C, D \in \mathbb{F}_{2^{n}}$, that the linear fractional transformation

$$
\mu(X)=\frac{A X+B}{C X+D},
$$

permutes $\mathbb{P}^{1}\left(\mathbb{F}_{2^{n}}\right)$ with the convention that $\mu(\infty)=A / C$. A projective polynomial over $\mathbb{F}_{2^{n}}$ is a polynomial

$$
A X^{q+1}+B X^{q}+C X+D, \quad A, B, C, D \in \mathbb{F}_{2^{n}}
$$

where $q=2^{i}$ for an integer $i$. Abhyankar started the work on these polynomials [10]. Bluher's paper [11] has become a standard reference for its detailed analysis of the roots of projective polynomials. We refer to [12] for more on Boolean functions and to [13] for finite fields.

\section*{III. Biprojective ApN Functions}
Let $F: \mathbb{F}_{2}^{n} \rightarrow \mathbb{F}_{2}^{n}$ be a vectorial Boolean function. In this section, we will assume that $n=2 m$. We can write $\mathbb{F}_{2^{2 m}}=\mathbb{F}_{2^{m}}+\beta \mathbb{F}_{2^{m}}$ where $\beta \in \mathbb{F}_{2^{2 m}} \backslash \mathbb{F}_{2^{m}}$. We will identify $\mathbb{F}_{2^{m}} \times \mathbb{F}_{2^{m}}$ with $\mathbb{F}_{2^{m}}+\beta \mathbb{F}_{2^{m}}$. Thus, the function $F: \mathbb{F}_{2}^{n} \rightarrow \mathbb{F}_{2}^{n}$ can be represented as a (univariate) polynomial (function) $F^{\prime}(X) \in \mathbb{F}_{2^{2 m}}[X]$, or a (bivariate) polynomial (function) $F^{\prime \prime}(x, y)=[f(x, y), g(x, y)] \in \mathbb{F}_{2^{m}}[x, y] \times \mathbb{F}_{2^{m}}[x, y]$. We will make no distinction between functions and the uni- or bivariate polynomials representing them and refer to all of them as $F .^{1}$

Now, as an example, consider the function $\kappa^{\prime}$ on $\mathbb{F}_{2^{6}}$, which is CCZ-equivalent to the Kim function $\kappa$ :

$$
\kappa^{\prime}(X)=X^{3}+X^{10}+b X^{17}
$$

\footnotetext{${ }^{1}$ The functions denoted by $F^{\prime}$ and $F^{\prime \prime}$ in this paper are not the derivatives of $F$.
}
where $b \in \mathbb{F}_{2^{3}}$ is a root of $x^{3}+x+1$. We can write any $X \in \mathbb{F}_{2^{6}}$ as
$$
X=x+\omega y
$$
where $\omega \in \mathbb{F}_{4}$ is a root of $z^{2}+z+1$ and $x, y \in \mathbb{F}_{2^{3}}$. Note that $\omega^{3}=\omega+\omega^{2}=1$. Thus we can write,
$$
\begin{aligned}
\kappa^{\prime}(X) & =X^{3}+X^{10}+b X^{17} \\
& =(x+\omega y)^{3}+(x+\omega y)^{10}+b(x+\omega y)^{17} \\
& =x^{3}+\omega x^{2} y+(\omega+1) x y^{2}+y^{3} \\
& +x^{3}+(\omega+1) x^{2} y+(\omega+1) x y^{2}+\omega y^{3} \\
& +b x^{3}+b \omega x^{2} y+b \omega x y^{2}+b(\omega+1) y^{3} \\
& =\omega\left(b x^{2} y+b x y^{2}+(b+1) y^{3}\right) \\
& +b x^{3}+x^{2} y+(b+1) y^{3}=\kappa^{\prime}(x, y) .
\end{aligned}
$$

We see that both $b x^{2} y+b x y^{2}+(b+1) y^{3}$ and $b x^{3}+ x^{2} y+(b+1) y^{3}$ are biprojective polynomials, i.e., projective polynomials in either variable $x, y$, when the other variable is fixed. This form is generalized in the following definition.

Definition III.1: Let


\begin{equation*}
F(x, y)=[f(x, y), g(x, y)] \in \mathbb{F}_{2^{m}}[x, y] \times \mathbb{F}_{2^{m}}[x, y], \tag{2}
\end{equation*}


with $q=2^{i}, r=2^{j}, i, j \geq 0$, and

$$
\begin{aligned}
& f(x, y)=a_{0} x^{q+1}+b_{0} x^{q} y+c_{0} x y^{q}+d_{0} y^{q+1} \\
& g(x, y)=a_{1} x^{r+1}+b_{1} x^{r} y+c_{1} x y^{r}+d_{1} y^{r+1}
\end{aligned}
$$

for some $a_{0}, a_{1}, b_{0}, b_{1}, c_{0}, c_{1}, d_{0}, d_{1} \in \mathbb{F}_{2^{m}}$. Then,\\
i) $f(x, y)$ is said to be a $q$-biprojective polynomial and $F(x, y)$ a $(q, r)$-biprojective polynomial pair.\\
ii) $f(x, y)$ is called the left-part of $F(x, y)$, and $g(x, y)$ is called the right-part of $F(x, y)$.\\
The following shorthand notation is used:

$$
\begin{aligned}
& a_{0} x^{q+1}+b_{0} x^{q} y+c_{0} x y^{q}+d_{0} y^{q+1}=\left(a_{0}, b_{0}, c_{0}, d_{0}\right)_{q}, \\
& a_{1} x^{r+1}+b_{1} x^{r} y+c_{1} x y^{r}+d_{1} y^{r+1}=\left(a_{1}, b_{1}, c_{1}, d_{1}\right)_{r} .
\end{aligned}
$$

\section*{A. A Survey on Known Biprojective APN Functions}
Now we will show that many known quadratic APN families and functions fall into the biprojective setting.

\begin{itemize}
  \item We have already shown that $\kappa^{\prime}(X)$ (also $\kappa(X)$ whose representation is slightly more complicated) can be represented as a ( 2,2 )-biprojective polynomial pair, since:
\end{itemize}

$$
\kappa^{\prime}(x, y)=\left[(b, 1,0, b+1)_{2},(0, b, b, b+1)_{2}\right]
$$

where $b$ is a root of $x^{3}+x+1$.

\begin{itemize}
  \item It is clear that the Gold functions $\mathrm{G}_{i}(X)=X^{2^{i}+1}$ can be written as ( $q, q$ )-biprojective polynomial pairs, as\\
$(x+\beta y)^{2^{i}+1}=x^{2^{i}+1}+\beta\left(x^{2^{i}} y\right)+\beta^{2^{i}}\left(x y^{2^{i}}\right)+\beta^{2^{i}+1} y^{2^{i}+1}$ shows. When $m$ is odd, we can use $\beta=\omega \in \mathbb{F}_{4} \backslash \mathbb{F}_{2}$. In that case\\
$(\mathcal{G}) \mathrm{G}_{i}(x, y)=\left[(1,0,1,1)_{2^{i}},(0,1,1,0)_{2^{i}}\right]$.
  \item Zhou and Pott found the following APN family [14] which also falls into the scheme of ( $q, r$ )-biprojective polynomial pairs. The functions\\
$(\mathcal{Z} \mathcal{P}) \quad F(x, y)=\left[(1,0,0, d)_{2^{i}},(0,0,1,0)_{2^{j}}\right], \quad d \in \mathbb{F}_{2^{m}}^{\times}$,\\
are APN if and only if $\operatorname{gcd}(i, m)=1, m$ is even and $d \neq a^{2^{i}+1}\left(b^{2^{i}}+b\right)^{1-2^{j}}$ for any $a, b \in \mathbb{F}_{2^{m}}$.
  \item In [15], Taniguchi found the following APN family, which is of the form (2):\\
(T) $F(x, y)=\left[(1,0, c, d)_{2^{i}},(0,0,1,0)_{2^{2 i}}\right], \quad c, d \in \mathbb{F}_{2^{m}}^{\times}$, where $\operatorname{gcd}(i, m)=1, f(x, 1) \neq 0$ for any $x \in \mathbb{F}_{2^{m}}$. If $c=0$ and $m$ is even, then the corresponding function belongs to the Zhou-Pott family [15].
  \item When $i=0$ (or $j=0$ ) in (2), i.e., $q=2^{0}=1$ and $r=2^{j}$ with $j>0$, then the resulting biprojective function can be written up to equivalence\\
(C) $F(x, y)=\left[x y,(a, b, c, d)_{r}\right]$,\\
which is the family $\mathcal{C}$ introduced by Carlet [16]. Indeed, Carlet showed that $F$ is APN if and only if $a \neq 0$ and $g(x, 1) \neq 0$ for any $x \in \mathbb{F}_{2^{m}}$.
\end{itemize}

\section*{B. The New Biprojective APN Functions}
In this paper we present two infinite families of ( $q, r$ )-biprojective polynomial pairs that are APN functions.

Theorem III.2: The following ( $q, r$ )-biprojective polynomial pairs $F(x, y)=[f(x, y), g(x, y)]$ are APN on $\mathbb{F}_{2^{m}} \times \mathbb{F}_{2^{m}}$.\\
$\left(\mathcal{F}_{1}\right)$ If $\operatorname{gcd}(3 i, m)=1$, then ${ }^{2}$

$$
F=\left[(1,0,1,1)_{2^{i}},(1,1,0,1)_{2^{2 i}}\right]
$$

is APN .\\
$\left(\mathcal{F}_{2}\right)$ If $\operatorname{gcd}(3 i, m)=1, m$ odd, then

$$
F=\left[(1,0,1,1)_{2^{i}},(0,1,1,0)_{2^{3 i}}\right]
$$

is APN .\\
$\left(\mathcal{F}_{3}\right)$ If $m=5$, then

$$
F=\left[(0,1,1,0)_{2},\left(1, u, u^{3}, u^{13}\right)_{2^{3}}\right]
$$

where $u$ is a root of $x^{5}+x^{2}+1$, is APN.\\
Remark III.3: Note that the family $\mathcal{F}_{2}$ can be seen as a Gold-hybrid APN function, i.e., a mix of the left-part $f(x, y)$ of $\mathrm{G}_{i}$ and the right-part $g(x, y)$ of $\mathrm{G}_{3 i}$. Note also that $F= \left[(1,0,1,1)_{2^{i}},(1,1,0,1)_{2^{2 i}}\right]$ in family $\mathcal{F}_{1}$ is EA-equivalent to $F^{\prime}=\left[(1,0,1,1)_{2^{i}},(1,0,1,1)_{2^{m-2 i}}\right]$ by $F^{\prime}=L \circ F$, where $L(x, y)=\left(x, y^{2^{m-2 i}}\right)$. Now when $m$ is odd, $F^{\prime}$ is also a Gold-hybrid APN function as the mix of the left-parts of $\mathrm{G}_{i}$ and $\mathrm{G}_{m-2 i}$. When $m$ is even, such a property does not seem to hold for family $\mathcal{F}_{1}$.

In [1], a necessary and sufficient condition was given for an APN function $F$ to be CCZ-equivalent to a permutation. We define

$$
Z_{\widehat{F}}=\left\{(U, V) \in \mathbb{F}_{2^{n}} \times \mathbb{F}_{2^{n}} \widehat{F}(U, V)=0\right\} \cup\{(0,0)\} .
$$

The function $F$ is CCZ-equivalent to a permutation if and only if $Z_{\widehat{F}}$ contains two $n$-dimensional, trivially intersecting $\mathbb{F}_{2}$ vector spaces [1]. Now assume $n=2 m$. In [4], a necessary condition for an APN function $F$ to be CCZ-equivalent to

\footnotetext{${ }^{2}$ Budaghyhan et al. [7], using a univariate approach, found APN families that intersect our family $\mathcal{F}_{1}$. The family $\mathcal{F}_{1}$ of our paper works for any $(i, 2 i)$ with $\operatorname{gcd}(i, m)=1$, while the functions of Budaghyan et al. that coincide with the family $\mathcal{F}_{1}$ correspond to the cases $i \in\{ \pm 1, \pm 1 / 2\}$ (the set is to be considered modulo $m$ ). For $m \geq 7$, our family $\mathcal{F}_{1}$ contains different functions as we will show in Section VI. Moreover, they require $m$ to be odd which we do not require. Both of the results require $\operatorname{gcd}(3, m)=1$. The methods and scopes of our respective papers are very different. See [7] for details.
}
a permutation is deduced from this observation: if $F$ is CCZ-equivalent to a permutation, then $\eta_{m}\left(\mathrm{NB}_{F}\right) \geq 2$. The only quadratic APN functions that satisfy this condition to the best of our knowledge are the Gold functions when $m$ is odd and the sporadic Kim function when $m=3$ (see [4] for computational results for small $n$ ). The function $\mathcal{F}_{3}$ (defined when $m=5$ ) satisfies this condition and seems to be sporadic like the Kim function $\kappa$. Even the condition $\eta_{m}\left(\mathrm{NB}_{F}\right) \geq 1$ seems to be very strong. Apart from the mentioned functions, the only known family (to the author) that satisfies this condition seems to be $X^{3}+A^{-1} \operatorname{tr}\left(A^{3} X^{9}\right)$ (found in [17], [18]) when $m$ is odd. The functions from family $\mathcal{F}_{2}$ of this paper satisfy this strong property (i.e., $\eta_{m}\left(\mathrm{NB}_{F}\right) \geq 1$ ). We summarize this discussion with the following proposition.

Proposition III.4: We have the following bounds for the following APN functions $F: \mathbb{F}_{2^{2 m}} \rightarrow \mathbb{F}_{2^{2 m}}$, when $m$ is odd.\\
i) If $F(X)=X^{2^{k}+1}$ with $\operatorname{gcd}(k, n)=1$, then $\eta_{m}\left(\mathrm{NB}_{F}\right)=\left(2^{m}+1\right) / 3$.\\
ii) if $F=\kappa$, then $\eta_{m}\left(\mathrm{NB}_{F}\right)=2$.\\
iii) if $F \in \mathcal{F}_{3}$, then $\eta_{m}\left(\mathrm{NB}_{F}\right)=2$.\\
iv) If $F(X)=X^{3}+\operatorname{tr}\left(X^{9}\right)$, then $\eta_{m}\left(\mathrm{NB}_{F}\right) \geq 1$.\\
v) If $F \in \mathcal{F}_{2}$, then $\eta_{m}\left(\mathrm{NB}_{F}\right) \geq 1$.

Proof: i) By Carlitz' Lemma ( [19], see [4, Lemma 1]) $\widehat{F}(A, 0) \notin\left\{ \pm 2^{m}\right\}$ if and only if $A$ is a nonzero cube in $\mathbb{F}_{2^{2 m}}$. When $m$ is odd, $3 \mid 2^{m}+1$ and all elements of $\mathbb{F}_{2^{m}}^{\times}$are cubes. Thus whenever $\alpha$ is a nonzero cube, $\alpha \mathbb{F}_{2^{m}}$ give $\left(2^{m}+1\right) / 3$ distinct $m$-dimensional subspaces consisting only of cubes (and zero). An $m$-dimensional subspace in $\mathbb{F}_{2^{2 m}}$ that consists only of cubes (and zero) is necessarily of the form $\alpha \mathbb{F}_{2^{m}}$ by [20, Theorem 1.1], which proves the equality.\\
ii) By computer.\\
iii) By computer.\\
iv) If $\operatorname{tr}(A)=0$, then $\widehat{F}(A, 0)=\widehat{\mathrm{G}_{1}}(A, 0)$ where $\mathrm{G}_{1}(X)$ is the Gold function $X \mapsto X^{3}$, since $\operatorname{tr}(X \operatorname{tr}(Y))= \operatorname{tr}(\operatorname{tr}(X) Y)$ for all $x, y \in \mathbb{F}_{2^{n}}$. By Part (i) and the fact that $H \cap \mathbb{F}_{2^{m}}=\mathbb{F}_{2^{m}}$, we get the result.\\
v) Recall that $q=2^{3 i}$ with $\operatorname{gcd}(3 i, m)=1$ and $m$ odd. For $a \in \mathbb{F}_{2^{m}}^{\times}$,

$$
\begin{aligned}
\widehat{F}((0, a),(0,0)) & =\sum_{x, y \in \mathbb{F}_{2^{m}}}(-1)^{\operatorname{tr}\left(a\left(x^{q} y+x y^{q}\right)\right)} \\
& =2^{m}+\sum_{\substack{x \in \mathbb{F}_{2^{m}} \\
y \in \mathbb{F}_{2^{m}}^{\infty}}}(-1)^{\operatorname{tr}\left(a y^{q+1}\left(x^{q}+x\right)\right)} \\
& =\sum_{\substack{x \in \mathbb{F}_{2^{m}} \\
y \in \mathbb{F}_{2} m}}(-1)^{\operatorname{tr}\left(a y\left(x^{q}+x\right)\right)} \\
& =\sum_{x \in \mathbb{F}_{2} m} \sum_{y \in \mathbb{F}_{2} m}(-1)^{\operatorname{tr}\left(y\left(x^{q}+x\right)\right)} \\
& =2^{m+1}
\end{aligned}
$$

since $x \mapsto x^{q+1}$ is a permutation of $\mathbb{F}_{2^{m}}$ and $\sum_{y \in \mathbb{F}_{2} m}(-1)^{\operatorname{tr}(y z)}$ is equal to $2^{m}$ if and only if $z=0$ and vanishes otherwise. After stripping $y=0$, we substituted $x y$ for $x$ in the second step and included $y=0$ back in the third step. The result follows from the condition $\operatorname{gcd}(3 i, m)=1$ since it implies $x^{q}+x=0$ has exactly two solutions.

Note that all previously known biprojective APN functions (other than the Gold functions) have one part (equivalent to) $(0,0,1,0)_{q}$. This corresponds (up to EA-equivalence) to the $(n, m)$-bent function $(x, y) \mapsto x y$ (i.e., all of its $2^{m}-1$ nonzero component functions are Boolean bent functions). Our examples are interesting in this respect as none of their parts are equivalent to $x y$. Although choosing a complicated left- (or right-) part that does not contain a bent component makes the proofs rather difficult, once done it leads to $\eta_{m}\left(\mathrm{NB}_{F}\right) \geq 1$ automatically as in the case of family $\mathcal{F}_{2}$. An interesting question is whether one can combine two parts that do not contain any bent component (not necessarily biprojective) to get APN functions with $\eta_{m}\left(\mathrm{NB}_{F}\right) \geq 2$. This function would then satisfy the necessary condition for being CCZ-equivalent to a permutation.

Problem III.5: Find a quadratic APN function F (or a family of functions) that satisfies $\eta_{m}\left(\mathrm{NB}_{F}\right) \geq 2$, possibly by generalizing $\mathcal{F}_{3}$.

We were not able to find other biprojective APN functions for small odd values of $m$ in our extensive search for such functions when $m \in\{3,5,7\}$. Details of these observations and negative results on biprojective APN functions will be the topic of a forthcoming companion paper.

\section*{C. A Projective APN Criterion}
The following is a simplification of the APN criterion for biprojective functions. For any ( $q, r$ )-biprojective polynomial $F(x, y)$ we define the derivatives

$$
\begin{aligned}
D_{f}^{[v, w]}(x, y) & =f(x+v, y+w)+f(x, y)+f(v, w) \\
& =\left(a_{0} v+b_{0} w\right) x^{q}+\left(a_{0} v^{q}+c_{0} w^{q}\right) x \\
& +\left(c_{0} v+d_{0} w\right) y^{q}+\left(b_{0} v^{q}+d_{0} w^{q}\right) y,
\end{aligned}
$$

and

$$
\begin{aligned}
D_{g}^{[v, w]}(x, y) & =g(x+v, y+w)+g(x, y)+g(v, w) \\
& =\left(a_{1} v+b_{1} w\right) x^{r}+\left(a_{1} v^{r}+c_{1} w^{r}\right) x \\
& +\left(c_{1} v+d_{1} w\right) y^{r}+\left(b_{1} v^{r}+d_{1} w^{r}\right) y .
\end{aligned}
$$

Now, $F$ is APN if and only if for any $(v, w) \neq(0,0)$,

$$
D_{f}^{[v, w]}(x, y)=0=D_{g}^{[v, w]}(x, y)
$$

has exactly two solutions

$$
(x, y) \in\{(0,0),(v, w)\} .
$$

Now, we will show that concentrating on the projective variant $\mathrm{D}^{u}$ of $D^{[v, w]}$ with respect to $u=v / w$, that is to say, $u \in \mathbb{P}^{1}\left(\mathbb{F}_{2^{m}}\right)$ is enough, because of the projective nature of our polynomials. If $v=0$ and $w \in \mathbb{F}_{2^{m}}^{\times}$, then we define,

$$
\begin{aligned}
& \mathrm{D}_{f}^{0}(x, y)=b_{0} x^{q}+c_{0} x+d_{0} y^{q}+d_{0} y=\frac{D_{f}^{[0, w]}(x w, y w)}{w^{q+1}} \\
& \mathrm{D}_{g}^{0}(x, y)=b_{1} x^{r}+c_{1} x+d_{1} y^{r}+d_{1} y=\frac{D_{g}^{[0, w]}(x w, y w)}{w^{r+1}} .
\end{aligned}
$$

Similarly, if $v \in \mathbb{F}_{2^{m}}^{\times}$and $w=0$, then we define,

$$
\begin{aligned}
& \mathrm{D}_{f}^{\infty}(x, y)=a_{0} x^{q}+a_{0} x+c_{0} y^{q}+b_{0} y=\frac{D_{f}^{[v, 0]}(x v, y v)}{v^{q+1}} \\
& \mathrm{D}_{g}^{\infty}(x, y)=a_{1} x^{r}+a_{1} x+c_{1} y^{r}+b_{1} y=\frac{D_{g}^{[v, 0]}(x v, y v)}{v^{r+1}}
\end{aligned}
$$

Now if $v w \neq 0$, then for $u=v / w$, define,

$$
\begin{aligned}
\mathrm{D}_{f}^{u}(x, y) & =\left(a_{0} u+b_{0}\right) x^{q}+\left(a_{0} u^{q}+c_{0}\right) x \\
& +\left(c_{0} u+d_{0}\right) y^{q}+\left(b_{0} u^{q}+d_{0}\right) y \\
& =\frac{D_{f}^{[v, w]}(x w, y w)}{w^{q+1}},
\end{aligned}
$$

and similarly

$$
\begin{aligned}
\mathrm{D}_{g}^{u}(x, y) & =\left(a_{1} u+b_{1}\right) x^{r}+\left(a_{1} u^{r}+c_{1}\right) x \\
& +\left(c_{1} u+d_{1}\right) y^{r}+\left(b_{1} u^{r}+d_{1}\right) y \\
& =\frac{D_{g}^{[v, w]}(x w, y w)}{w^{r+1}}
\end{aligned}
$$

The following lemma summarizes the above discussion.\\
Lemma III.6: Let $F(x, y)=[f(x, y), g(x, y)]$ be a $(q, r)$ biprojective polynomial pair. Then $F$ is APN if and only if $\mathrm{D}_{f}^{u}(x, y)=0=\mathrm{D}_{g}^{u}(x, y)$ has exactly two solutions for each $u \in \mathbb{P}^{1}\left(\mathbb{F}_{2^{m}}\right)$.

\section*{IV. The Method}
For $u \in \mathbb{P}^{1}\left(\mathbb{F}_{2^{m}}\right)$ we will analyze the solutions of $\mathrm{E}_{f}^{u}(x, y)=0$ and $\mathrm{E}_{g}^{u}(x, y)=0$ defined below, which are slight variations of $\mathrm{D}_{f}^{u}$ and $\mathrm{D}_{g}^{u}$.

\begin{itemize}
  \item For $u \in\{0, \infty\}$, define $\mathrm{E}_{f}^{u}(x, y)=\mathrm{D}_{f}^{u}(x, y)$, and $\mathrm{E}_{g}^{u}(x, y)=\mathrm{D}_{g}^{u}(x, y)$.
  \item For $u \in \mathbb{F}_{2^{m}}^{\times}$, by substituting $x \mapsto x u$ and then $y \mapsto x+y$, define
\end{itemize}

$$
\begin{aligned}
\mathrm{E}_{f}^{u}(x, y) & =f(u, 1)\left(x^{q}+x\right) \\
& +\left(c_{0} u+d_{0}\right) y^{q}+\left(b_{0} u^{q}+d_{0}\right) y \\
& =\mathrm{D}_{f}^{u}(x u, y+x) \\
\mathrm{E}_{g}^{u}(x, y) & =g(u, 1)\left(x^{r}+x\right) \\
& +\left(c_{1} u+d_{1}\right) y^{r}+\left(b_{1} u^{r}+d_{1}\right) y \\
& =\mathrm{D}_{g}^{u}(x u, y+x)
\end{aligned}
$$

The pair of equations $\mathrm{E}_{f}^{u}(x, y)=0=\mathrm{E}_{g}^{u}(x, y)$ has the same number of common solutions as the pair of equations $\mathrm{D}_{f}^{u}(x, y)=0=\mathrm{D}_{g}^{u}(x, y)$. We will use the following variant of Lemma III.6.

Lemma IV.1: Let $F(x, y)=[f(x, y), g(x, y)]$ be a $(q, r)$ biprojective polynomial pair. Then $F$ is APN if and only if $\mathrm{E}_{f}^{u}(x, y)=0=\mathrm{E}_{g}^{u}(x, y)$, has exactly two solutions for each $u \in \mathbb{P}^{1}\left(\mathbb{F}_{2^{m}}\right)$.

\section*{A. A Linear Transformation Lemma}
Let

$$
\begin{aligned}
F(x, y) & =\left[\left(a_{0}, b_{0}, c_{0}, d_{0}\right)_{q},\left(a_{1}, b_{1}, c_{1}, d_{1}\right)_{r}\right] \\
& =[f(x, y), g(x, y)]
\end{aligned}
$$

with $q=2^{i}$ and $r=2^{j}$ with $1 \leq i, j \leq m-1$. Assume that $\operatorname{gcd}(i, m)=1$. Then, there exists $1 \leq k \leq m-1$ such that $i k \equiv j(\bmod m)$, and we have $r \equiv q^{k}\left(\bmod 2^{m}-1\right)$. Thus, we can re-write the equations $\mathrm{E}_{f}^{u}(x, y)=0=\mathrm{E}_{g}^{u}(x, y)$ with $u \in \mathbb{F}_{2^{m}}^{\times}$as


\begin{align*}
\psi_{q}(x) & =M_{f}^{u}(y),  \tag{3}\\
\psi_{q^{k}}(x) & =M_{g}^{u}(y), \tag{4}
\end{align*}


where

$$
\begin{aligned}
\psi_{2^{i}}(x) & =x^{2^{i}}+x \\
M_{f}^{u}(y) & =\frac{\left(c_{0} u+d_{0}\right) y^{q}+\left(b_{0} u^{q}+d_{0}\right) y}{f(u, 1)} \\
M_{g}^{u}(y) & =\frac{\left(c_{1} u+d_{1}\right) y^{q^{k}}+\left(b_{1} u^{q^{k}}+d_{1}\right) y}{g(u, 1)}
\end{aligned}
$$

whenever $f(u, 1) g(u, 1) \neq 0$. Consider the $\mathbb{F}_{2}$-linear operator

$$
\mathcal{L}(\alpha, \beta)=\sum_{l=0}^{k-1} \alpha^{q^{l}}+\beta
$$

It is easy to see

$$
\mathcal{L}\left(\psi_{q}(x), \psi_{q^{k}}(x)\right)=0
$$

Define

$$
L_{F}^{u}(y)=\mathcal{L}\left(M_{f}^{u}(y), M_{g}^{u}(y)\right)
$$

We thus have for all $x, y \in \mathbb{F}_{2^{m}}$,


\begin{equation*}
\mathcal{L}\left(\psi_{q}(x)+M_{f}^{u}(y), \psi_{q^{k}}(x)+M_{g}^{u}(y)\right)=L_{F}^{u}(y) \tag{5}
\end{equation*}


Now, if $\left(x_{0}, y_{0}\right)$ is a common solution of (3) and (4), then

$$
\psi_{q}\left(x_{0}\right)+M_{f}^{u}\left(y_{0}\right)=0=\psi_{q^{k}}\left(x_{0}\right)+M_{g}^{u}\left(y_{0}\right)
$$

therefore,

$$
\mathcal{L}\left(\psi_{q}\left(x_{0}\right)+M_{f}^{u}\left(y_{0}\right), \psi_{q^{k}}\left(x_{0}\right)+M_{g}^{u}\left(y_{0}\right)\right)=L_{F}^{u}\left(y_{0}\right)=0 .
$$

Hence, $\operatorname{Ker}\left(L_{F}^{u}\right)$ contains all $y \in \mathbb{F}_{2^{m}}$ such that $(x, y)$ is a common solution of $\mathrm{E}_{f}^{u}(x, y)=0=\mathrm{E}_{g}^{u}(x, y)$ for some $x \in \mathbb{F}_{2^{m}}$. Now assume that the pair of equations $\mathrm{E}_{f}^{u}(x, y)= 0=\mathrm{E}_{g}^{u}(x, y)$ has exactly two common solutions for each $u$ in

$$
\{0, \infty\} \cup\left\{u \in \mathbb{F}_{2^{m}}^{\times} \mid f(u, 1) g(u, 1)=0\right\} .
$$

Since $\operatorname{gcd}(i, m)=1$, we have $\operatorname{Ker}\left(\psi_{q}\right) \cap \operatorname{Ker}\left(\psi_{q^{k}}\right)=\mathbb{F}_{2}$. Now, $F$ is APN if and only if for all the remaining $u \in \mathbb{F}_{2^{m}}^{\times}$(i.e., $f(u, 1) g(u, 1) \neq 0$ ), we have either,\\
i) $\operatorname{Ker}\left(L_{F}^{u}\right)=\{0\}$; or,\\
ii) $\operatorname{Ker}\left(L_{F}^{u}\right) \supsetneq\{0\}$ and for all nonzero $y_{u} \in \operatorname{Ker}\left(L_{F}^{u}\right) \backslash\{0\}$, we have


\begin{equation*}
\psi_{q}(x) \neq M_{f}^{u}\left(y_{u}\right) \tag{6}
\end{equation*}


or


\begin{equation*}
\psi_{q^{k}}(x) \neq M_{g}^{u}\left(y_{u}\right), \tag{7}
\end{equation*}


for every $x \in \mathbb{F}_{2^{m}}$.\\
Indeed, for a given $u$, the common solutions of (3) and (4) are precisely $(x, y) \in \mathbb{F}_{2} \times\{0\}$ if and only if either $\operatorname{Ker}\left(L_{F}^{u}\right)= \{0\}$, or $\operatorname{Ker}\left(L_{F}^{u}\right) \neq\{0\}$ and at least one of the conditions (6) and (7) holds for every $y_{u} \in \operatorname{Ker}\left(L_{F}^{u}\right) \backslash\{0\}$ and $x \in \mathbb{F}_{2^{m}}$.

Since $\operatorname{gcd}(i, m)=1$ and therefore $\psi_{q}(x)$ is 2-to-1 onto the trace-0 hyperplane of $\mathbb{F}_{2^{m}}$, the condition (6) is equivalent to $\operatorname{tr}\left(M_{f}^{u}\left(y_{u}\right)\right)=1$, for all $u \in \mathbb{F}_{2^{m}}^{\times}$with $f(u, 1) g(u, 1) \neq 0$ and $y_{u} \in \operatorname{Ker}\left(L_{F}^{u}\right) \backslash\{0\}$. Note that for $y_{u} \in \operatorname{Ker}\left(L_{F}^{u}\right)$ the condition (7) implies the condition (6) by (5) and the definition of $\mathcal{L}$.

We record the above observation as a lemma.\\
Lemma IV.2: Let

$$
\begin{aligned}
F(x, y) & =\left[\left(a_{0}, b_{0}, c_{0}, d_{0}\right)_{q},\left(a_{1}, b_{1}, c_{1}, d_{1}\right)_{q^{k}}\right] \\
& =[f(x, y), g(x, y)]
\end{aligned}
$$

with $1 \leq i, k \leq m-1, q=2^{i}$ and $\operatorname{gcd}(i, m)=1$. Let $M_{f}^{u}$ and $L_{F}^{u}$ be defined as above. Then, $F$ is APN if and only if

\begin{itemize}
  \item the pair of equations $\mathrm{E}_{f}^{u}(x, y)=0=\mathrm{E}_{g}^{u}(x, y)$ has exactly two common solutions for $u \in\{0, \infty\} \cup\{u \in \left.\mathbb{F}_{2^{m}}^{\times} \mid f(u, 1) g(u, 1)=0\right\}$, and
  \item for all $u \in \mathbb{F}_{2^{m}}^{\times}$with $f(u, 1) g(u, 1) \neq 0$ and $y_{u} \in \operatorname{Ker}\left(L_{F}^{u}\right) \backslash\{0\}$, we have $\operatorname{tr}\left(M_{f}^{u}\left(y_{u}\right)\right)=1$.\\
Remark IV.3: If $\operatorname{Ker}\left(L_{F}^{u}\right)=\{0\}$ for some $u$, then the trace condition is vacuously true.
\end{itemize}

In the next sections, we are going to show that for our families the corresponding $L_{F}^{u}(y)$ defined above are all related to each other projectively. For that purpose we will use a projective transformation that keeps the number of zeroes of a $q$-projective polynomial invariant (see Lemma IV. 4 (iii)).

\section*{B. Projective Polynomials}
Note that $f(u, 1)$ and $g(u, 1)$ are (univariate) projective polynomials. The following specific projective polynomial

$$
\phi_{q}(x)=x^{q+1}+x+1,
$$

is pivotal in this paper. The following lemma is a collection of rather straightforward results on $\phi$ and related polynomials required for the sequel.

Lemma IV.4: Let $\operatorname{gcd}(3 k, m)=1$ and $q=2^{k}$. Then\\
i) $\phi_{q}(x) \neq 0$ for $x \in \mathbb{F}_{2^{m}}$,\\
ii) $\phi_{q^{2}}(x)=0$ for some $x \in \mathbb{F}_{2^{m}}$ if and only if $m$ is even and $x \in \mathbb{F}_{4} \backslash \mathbb{F}_{2}$,\\
iii) if $\epsilon_{1} \epsilon_{4}+\epsilon_{3} \epsilon_{2} \neq 0$ then for any $x \in \mathbb{F}_{2^{m}}$, we have,

$$
\phi^{\prime}(x)=\left(\epsilon_{3} x+\epsilon_{4}\right)^{q+1} \phi_{q}\left(\frac{\epsilon_{1} x+\epsilon_{2}}{\epsilon_{3} x+\epsilon_{4}}\right) \neq 0
$$

where $\epsilon_{1}, \epsilon_{2}, \epsilon_{3}, \epsilon_{4} \in \mathbb{F}_{2^{m}}$,\\
iv) if $f(x)=a x^{q+1}+b x+c \neq 0$ for all $x \in \mathbb{F}_{2^{m}}$, then the linearized polynomial

$$
l(x)=x f\left(x^{q-1}\right)=a x^{q^{2}}+b x^{q}+c x
$$

permutes $\mathbb{F}_{2^{m}}$.\\
Proof: i) All $q^{2}$ roots $y_{i}$ of $y \phi_{q}\left(y^{q-1}\right)=y+y^{q}+ y^{q^{2}}=\operatorname{tr}_{\mathbb{F}_{q^{3}} / \mathbb{F}_{q}}(y)$ are in $\mathbb{F}_{q^{3}}$. All $q+1$ roots $z_{j}$ of $\phi_{q}(z)$ are therefore of the form $y_{i}^{q-1}=z_{j} \in \mathbb{F}_{q^{3}}$. But $\mathbb{F}_{2^{3 k}} \cap \mathbb{F}_{2^{m}}=\mathbb{F}_{2}$ because of the $\operatorname{gcd}$ assumption and consequently $\phi_{q}(z)$ has no roots in $\mathbb{F}_{2}$.\\
ii) Part (i) implies $m=2 m^{\prime}$ is even and $\mathbb{F}_{2^{6 k}} \cap \mathbb{F}_{2^{2 m^{\prime}}}=\mathbb{F}_{4}$. We have $2^{2 k} \equiv 1(\bmod 3)$ hence $\omega \in \mathbb{F}_{4} \backslash \mathbb{F}_{2}$ satisfies $\omega^{q^{2}+1}+\omega+1=\omega^{2}+\omega+1=0$ and the result follows.\\
iii) Note that $\phi^{\prime}(x)$ is a polynomial, i.e., $\left(\epsilon_{3} x+\epsilon_{4}\right)^{q+1}$ cancels the denominators appearing in the expansion of $\phi_{q}\left(\frac{\epsilon_{1} x+\epsilon_{2}}{\epsilon_{3} x+\epsilon_{4}}\right)$. The map

$$
x \mapsto\left(\frac{\epsilon_{1} x+\epsilon_{2}}{\epsilon_{3} x+\epsilon_{4}}\right)
$$

permutes $\mathbb{P}^{1}\left(\mathbb{F}_{2^{m}}\right)$. When the image is not $\infty$ for $x \in \mathbb{F}_{2^{m}}$, i.e., $x \neq \frac{\epsilon_{4}}{\epsilon_{3}}$, then $\phi^{\prime}(x)$ is a product of two $\mathbb{F}_{2^{n}}^{\times}$elements. Thus, the only nontrivial case is when $\infty$ appears, i.e., $x=\frac{\epsilon_{4}}{\epsilon_{3}}$ with $x \in \mathbb{F}_{2^{m}}$. For this case we have $\phi^{\prime}(x)=\left(\epsilon_{1} x+\epsilon_{2}\right)^{q+1} \neq 0$, since the determinant $\epsilon_{1} \epsilon_{4}+\epsilon_{3} \epsilon_{2}$ is nonzero and hence $\frac{e_{2}}{e_{1}} \neq \frac{e_{4}}{e_{3}}$.\\
iv) Obviously $\operatorname{Ker}(l)=\{0\}$.

We are now ready to prove the APN property of the families $\mathcal{F}_{1}, \mathcal{F}_{2}$ and the function $\mathcal{F}_{3}$.

\section*{V. Families $\mathcal{F}_{1}, \mathcal{F}_{2}$ and the Function $\mathcal{F}_{3}$}
In this section we prove Theorem III.2. Each subsection is devoted to a proof of a family.

\section*{A. Family $\mathcal{F}_{1}$}
Theorem V.1: Let $F(x, y)=[f(x, y), g(x, y)]$ be the vectorial Boolean function on $\mathbb{F}_{2^{m}} \times \mathbb{F}_{2^{m}}$ with

\begin{itemize}
  \item $\operatorname{gcd}(3 i, m)=1$ and $q=2^{i}$ with $1 \leq i \leq m-1$,
  \item $f=(1,0,1,1)_{q}$,
  \item $g=(1,1,0,1)_{q^{2}}$.
\end{itemize}

Then $F$ is APN.\\
Proof: First, from $\mathrm{E}_{f}^{0}(x, y)=0=\mathrm{E}_{g}^{0}(x, y)$, we get $x+y^{q}+y=0=x^{q^{2}}+y^{q^{2}}+y$ which implies $y+y^{q^{3}}=0$. Since $\operatorname{gcd}(3 i, m)=1, y \in\{0,1\}$ are the only solutions to $y+y^{q^{3}}=0$. Both of these solutions imply $x=0$, therefore $\mathrm{E}_{f}^{0}(x, y)=0=\mathrm{E}_{g}^{0}(x, y)$ has exactly two solutions: $(x, y) \in\{(0,1),(0,0)\}$. The case $u \in\{\infty, 1\}$ is identical up to a change of variables. Now assume $u \in \mathbb{F}_{2^{m}}^{\times}$. We have $\mathrm{E}_{f}^{u}(x, y)=0=\mathrm{E}_{g}^{u}(x, y)$, i.e.,

$$
\begin{aligned}
\left(u^{q+1}+u+1\right)\left(x^{q}+x\right)+(u+1) y^{q}+y & =0, \text { and } \\
\left(u^{q^{2}+1}+u^{q^{2}}+1\right)\left(x^{q^{2}}+x\right)+y^{q^{2}}+\left(u^{q^{2}}+1\right) y & =0 .
\end{aligned}
$$

According to Lemma A. 1 (i), we have exactly two common solutions when $u \in \mathbb{F}_{4} \backslash \mathbb{F}_{2}$. Therefore, we can concentrate on $u \in \mathbb{F}_{2^{m}} \backslash \mathbb{F}_{4}$. Since $\operatorname{gcd}(3 i, m)=1$, by Lemma IV. 4 (i) and (ii), we have

$$
\begin{aligned}
\phi_{q}(u) & =u^{q+1}+u+1 \neq 0, \\
\phi_{q^{2}}(u+1) & =u^{q^{2}+1}+u^{q^{2}}+1 \neq 0 .
\end{aligned}
$$

Therefore we can define two new equations:


\begin{align*}
x^{q}+x & =\mu_{u}(y), \text { and }  \tag{8}\\
x^{q^{2}}+x & =\nu_{u}(y), \tag{9}
\end{align*}


where

$$
\begin{aligned}
& \mu_{u}(y)=\frac{(u+1) y^{q}+y}{\phi_{q}(u)}, \text { and } \\
& \nu_{u}(y)=\frac{y^{q^{2}}+(u+1)^{q^{2}} y}{\phi_{q^{2}}(u+1)} .
\end{aligned}
$$

We will show that

$$
\lambda_{u}^{\prime \prime}(y)=\mu_{u}(y)+\mu_{u}(y)^{q}+\nu_{u}(y)
$$

is a permutation for every $u \in \mathbb{F}_{2^{m}} \backslash \mathbb{F}_{4}$. By definition


\begin{align*}
\lambda_{u}^{\prime \prime}(y) & =\left[\left(\frac{u+1}{\phi_{q}(u)}\right)^{q}+\frac{1}{\phi_{q^{2}}(u+1)}\right] y^{q^{2}} \\
& +\left[\frac{u+1}{\phi_{q}(u)}+\left(\frac{1}{\phi_{q}(u)}\right)^{q}\right] y^{q} \\
& +\left[\frac{(u+1)^{q^{2}}}{\phi_{q^{2}}(u+1)}+\frac{1}{\phi_{q}(u)}\right] y . \tag{10}
\end{align*}


Now define $\lambda_{u}^{\prime}(y)=\lambda_{u}^{\prime \prime}(y)\left(\phi_{q}(u)\right)^{q+1} \phi_{q^{2}}(u+1)$. We get by Lemma A. 1 (ii),


\begin{align*}
\lambda_{u}^{\prime}(y) & =u^{q^{2}}\left(\phi_{q}(u)\right)^{2} y^{q^{2}} \\
& +u^{q}\left(\phi_{q^{2}}(u+1)\right)^{2} y^{q} \\
& +u\left(\phi_{q}(u)\right)^{2 q} y \tag{11}
\end{align*}


Letting $\lambda_{u}(y)=\lambda_{u}^{\prime}(y / u)$,

$$
\lambda_{u}(y)=\left(\phi_{q}(u)\right)^{2} y^{q^{2}}+\left(\phi_{q^{2}}(u+1)\right)^{2} y^{q}+\left(\phi_{q}(u)\right)^{2 q} y .
$$

Now the projective polynomial defined by

$$
\pi_{u}\left(y^{q-1}\right)=\frac{\lambda_{u}(y)}{y}
$$

i.e.,

$$
\pi_{u}(x)=\left(\phi_{q}(u)\right)^{2} x^{q+1}+\left(\phi_{q^{2}}(u+1)\right)^{2} x+\left(\phi_{q}(u)\right)^{2 q}
$$

satisfies (by Lemma A. 1 (iii) and a simple application of a Galois automorphism)

$$
\pi_{u}(x)=\left(\epsilon_{3} x+\epsilon_{4}\right)^{q+1} \phi_{q}\left(\frac{\epsilon_{1} x+\epsilon_{2}}{\epsilon_{3} x+\epsilon_{4}}\right)
$$

with

$$
\left(\begin{array}{ll}
\epsilon_{1} & \epsilon_{2} \\
\epsilon_{3} & \epsilon_{4}
\end{array}\right)=\left(\begin{array}{cc}
1 & (u+1)^{2 q} \\
(u+1)^{2} & u^{2 q}
\end{array}\right),
$$

whose determinant is conveniently

$$
\left|\begin{array}{cc}
1 & (u+1)^{2 q} \\
(u+1)^{2} & u^{2 q}
\end{array}\right|=\left(\phi_{q}(u)\right)^{2} \neq 0
$$

for any $u \in \mathbb{F}_{2^{m}} \backslash \mathbb{F}_{4}$ by Lemma IV. 4 (i). Therefore by Lemma IV. 4 (iii) and (iv) we get the fact that the polynomial $\lambda_{u}(y)$ permutes $\mathbb{F}_{2^{m}}$, for all $u \in \mathbb{F}_{2^{m}} \backslash \mathbb{F}_{4}$.

Note that $\lambda_{u}^{\prime}(y)$ and in turn $\lambda_{u}^{\prime \prime}(y)$ are also permutations. Now, Lemma IV. 2 proves that the family is APN (setting $L_{F}^{u}=\lambda_{u}^{\prime \prime}, M_{f}^{u}=\mu_{u}$ and $M_{g}^{u}=\nu_{u}$ ).

We record the above observation as a lemma which we will also use for the next family.

Lemma V.2: For $u \in \mathbb{F}_{2^{m}} \backslash \mathbb{F}_{4}$, the polynomial $\lambda_{u}(y)$ permutes $\mathbb{F}_{2^{m}}$.

\section*{B. Family $\mathcal{F}_{2}$}
Theorem V.3: Let $F(x, y)=[f(x, y), g(x, y)]$ be the vectorial Boolean function on $\mathbb{F}_{2^{m}} \times \mathbb{F}_{2^{m}}$ with,

\begin{itemize}
  \item $m$ odd,
  \item $\operatorname{gcd}(3 i, m)=1$ and $q=2^{i}$ with $1 \leq i \leq m-1$,
  \item $f=(1,0,1,1)_{q}$,
  \item $g=(0,1,1,0)_{q^{3}}$.
\end{itemize}

Then $F$ is APN.\\
Proof: First, we consider

$$
\begin{aligned}
& \mathrm{E}_{f}^{0}(x, y)=x+y^{q}+y=0=x^{q^{3}}+x=\mathrm{E}_{g}^{0}(x, y), \text { and } \\
& \mathrm{E}_{f}^{\infty}(x, y)=x^{q}+x+y^{q}=0=y^{q^{3}}+y=\mathrm{E}_{g}^{\infty}(x, y) .
\end{aligned}
$$

The only solutions to these equation pairs are $(x, y) \in \{(0,0),(0,1)\}$ and $(x, y) \in\{(0,0),(1,0)\}$ respectively. Indeed, $\operatorname{gcd}(3 i, m)=1$ implies that the right-hand sides of the above equations are satisfied if and only if $x \in \mathbb{F}_{2}$ (resp. $y \in \mathbb{F}_{2}$ ). Now, the assumption that $m$ is odd implies that $\operatorname{tr}(1)=1$ and $y^{q}+y=1$ (resp. $x^{q}+x=1$ ) has no solutions.

If $u \in \mathbb{F}_{2^{m}}^{\times}$, then

$$
\begin{aligned}
& \mathrm{E}_{f}^{u}(x, y)=\left(u^{q+1}+u+1\right)\left(x^{q}+x\right)+(u+1) y^{q}+y=0 \\
& \mathrm{E}_{g}^{u}(x, y)=\left(u^{q^{3}}+u\right)\left(x^{q^{3}}+x\right)+u y^{q^{3}}+u^{q^{3}} y=0
\end{aligned}
$$

The case $u=1$ reduces to the case of $\mathrm{E}_{g}^{\infty}$ up to a change of variable. Now, assume $u \in \mathbb{F}_{2^{m}} \backslash \mathbb{F}_{2}$. By Lemma IV. 4 (i), we have $\phi_{q}(u)=u^{q+1}+u+1 \neq 0$, and by the assumption $\operatorname{gcd}(3 i, m)=1$, we have $\psi_{q^{3}}(u)=u^{q^{3}}+u \neq 0$ for $u \notin \mathbb{F}_{2}$. Now we want to count the common solutions of


\begin{align*}
x^{q}+x & =\mu_{u}(y), \text { and }  \tag{12}\\
x^{q^{3}}+x & =\sigma_{u}(y), \tag{13}
\end{align*}


where

$$
\begin{aligned}
& \mu_{u}(y)=\frac{(u+1) y^{q}+y}{\phi_{q}(u)}, \text { and } \\
& \sigma_{u}(y)=\frac{u y^{q^{3}}+u^{q^{3}} y}{u^{q^{3}}+u} .
\end{aligned}
$$

We are going to show that

$$
\tau_{u}^{\prime}(y)=\mu_{u}(y)+\mu_{u}(y)^{q}+\mu_{u}(y)^{q^{2}}+\sigma_{u}(y)
$$

is a 2 -to-1 map. By definition


\begin{align*}
\tau_{u}^{\prime}(y) & =\left[\left(\frac{u+1}{\phi_{q}(u)}\right)^{q^{2}}+\frac{u}{u^{q^{3}}+u}\right] y^{q^{3}} \\
& +\left[\left(\frac{u+1}{\phi_{q}(u)}\right)^{q}+\left(\frac{1}{\phi_{q}(u)}\right)^{q^{2}}\right] y^{q^{2}} \\
& +\left[\frac{u+1}{\phi_{q}(u)}+\left(\frac{1}{\phi_{q}(u)}\right)^{q}\right] y^{q} \\
& +\left[\frac{u^{q^{3}}}{u^{q^{3}}+u}+\frac{1}{\phi_{q}(u)}\right] y . \tag{14}
\end{align*}


We can show (see Lemma A. 2 (i))

$$
\begin{aligned}
\tau_{u}^{\prime}(y) & =\frac{u^{q^{3}} \phi_{q^{2}}(u+1)}{\left(\phi_{q}(u)\right)^{q^{2}} \psi_{q^{3}}(u)} y^{q^{3}}+\frac{u^{q^{2}}\left(\phi_{q^{2}}(u+1)\right)^{q}}{\left(\phi_{q}(u)\right)^{q^{2}+q}} y^{q^{2}} \\
& +\frac{u^{q} \phi_{q^{2}}(u+1)}{\left(\phi_{q}(u)\right)^{q+1}} y^{q}+\frac{u\left(\phi_{q^{2}}(u+1)\right)^{q}}{\phi_{q}(u) \psi_{q^{3}}(u)} y
\end{aligned}
$$

Now set

$$
\begin{aligned}
\tau_{u}(y) & =\tau_{u}^{\prime}\left(\frac{y}{u}\right) \frac{\left(\phi_{q}(u)\right)^{2 q+1} \psi_{q^{3}}(u)}{\left(\phi_{q^{2}}(u+1)\right)^{q}} \\
& =\frac{\left(\phi_{q}(u)\right)^{2 q}}{\left(\phi_{q^{2}}(u+1)\right)^{q-1}\left(\phi_{q}(u)\right)^{q^{2}-1}} y^{q^{3}} \\
& +\frac{\left(\phi_{q}(u)\right)^{q} \psi_{q^{3}}(u)}{\left(\phi_{q}(u)\right)^{q^{2}-1}} y^{q^{2}} \\
& +\frac{\left(\phi_{q}(u)\right)^{q} \psi_{q^{3}}(u)}{\left(\phi_{q^{2}}(u+1)\right)^{q-1}} y^{q}+\left(\phi_{q}(u)\right)^{2 q} y
\end{aligned}
$$

It can be shown that (see Lemma A. 2 (ii))

$$
\tau_{u}(y)=\lambda_{u}(y)+\frac{\left(\lambda_{u}(y)\right)^{q}}{C_{u}},
$$

where

$$
C_{u}=\left(\phi_{q^{2}}(u+1)\right)^{q-1}\left(\phi_{q}(u)\right)^{q^{2}-1},
$$

and

$$
\lambda_{u}(y)=\left(\phi_{q}(u)\right)^{2} y^{q^{2}}+\left(\phi_{q^{2}}(u+1)\right)^{2} y^{q}+\left(\phi_{q}(u)\right)^{2 q} y,
$$

which was defined for family $\mathcal{F}_{1}$ previously.

Claim: For $u \in \mathbb{F}_{2^{m}} \backslash \mathbb{F}_{2}$, the polynomial $\tau_{u}(y)$ is a 2-to-1 map with $\operatorname{Ker}\left(\tau_{u}\right)=\left\{0, y_{u}\right\}$, where $y_{u}=\lambda_{u}^{-1}\left(\phi_{q^{2}}(u+\right.$ 1) $\left.\left(\phi_{q}(u)\right)^{q+1}\right)$.

Proof of the Claim. Observe that $\lambda_{u}$ is a permutation and therefore $\operatorname{Ker}\left(\lambda_{u}\right)=\{0\}$. By Lemma IV. 4 (i) and (ii), $C_{u} \neq 0$. So any $y_{u} \in \operatorname{Ker}\left(\tau_{u}\right) \backslash\{0\}$ satisfies

$$
\lambda_{u}\left(y_{u}\right)^{q-1}=\left(\phi_{q^{2}}(u+1)\right)^{q-1}\left(\phi_{q}(u)\right)^{q^{2}-1}=C_{u} .
$$

Since $\operatorname{gcd}\left(q-1,2^{m}-1\right)=2^{\operatorname{gcd}(i, m)}-1=1$, the ( $q-1$ )st root of $C_{u}$ in $\mathbb{F}_{2^{m}}$ is unique. Therefore, $y_{u} \in \operatorname{Ker}\left(\tau_{u}\right) \backslash\{0\}$ which is necessarily mapped by $\lambda_{u}\left(y_{u}\right)^{q-1}$ to $C_{u}$ is unique. Thus, we have proved the claim that $\tau_{u}(y)$ is a 2 -to- 1 map with the given kernel.

Now it is clear that $\operatorname{Ker}\left(\tau_{u}^{\prime}\right)=\left\{0, y_{u} / u\right\}$. We easily show (see Lemma A. 2 (iii)) that $y_{u}=u^{2}+u+1$. Thus, $\frac{u^{2}+u+1}{u}$ is the only nonzero element in $\operatorname{Ker}\left(\tau_{u}^{\prime}\right)$, and we will show that $\operatorname{tr}\left(\mu_{u}\left(\frac{u^{2}+u+1}{u}\right)\right)=1$. Indeed, by Lemma A. 2 (iv),

$$
\mu_{u}\left(\frac{u^{2}+u+1}{u}\right)=\frac{1}{u^{q}}+\frac{1}{u}+1,
$$

thus the condition of Lemma IV. 2 is satisfied since $\operatorname{tr}\left(\frac{1}{u^{q}}+\frac{1}{u}+1\right)=1$ by our assumption that $m$ is odd. Therefore, by Lemma IV. 2 (setting $L_{F}^{u}=\tau_{u}^{\prime}, M_{f}^{u}=\mu_{u}$ and $M_{g}^{u}=\sigma_{u}$ ), we prove that $F$ is APN.

\section*{C. The Function $\mathcal{F}_{3}$}
Theorem V.4: Let $F(x, y)=[f(x, y), g(x, y)]$ be the vectorial Boolean function on $\mathbb{F}_{2^{5}} \times \mathbb{F}_{2^{5}}$ with,

\begin{itemize}
  \item $f=(0,1,1,0)_{2}$,
  \item $g=\left(1, u, u^{3}, u^{13}\right)_{2^{3}}$, where $u \in \mathbb{F}_{2^{5}}$ is a root of $x^{5}+x^{2}+1$.\\
Then $F$ is APN.\\
Proof: The fact that $F$ is APN can be verified easily with a computer.
\end{itemize}

\section*{VI. Inequivalence Results}
In this section we will show that the families $\mathcal{F}_{1}, \mathcal{F}_{2}$ and the function $\mathcal{F}_{3}$ give APN functions CCZ-inequivalent to functions from previously known families. But first, let us explain how the family $\mathcal{F}_{1}$ and the recent family $\mathcal{B} \mathcal{H} \mathcal{K}$ of Budaghyan et al. intersect. Recall that for two quadratic APN functions, EA-equivalence coincides with CCZ-equivalence by a result of Yoshiara [3].

\section*{A. Family $\mathcal{F}_{1}$ and the Family $\mathcal{B} \mathcal{H} \mathcal{K}$ of Budaghyan Et Al}
The following theorem was proved in [7].\\
Theorem VI.1: Let $n=2 m, m$ odd, $3 \nmid m$, and $\omega \in \mathbb{F}_{4}$ be a generator of $\mathbb{F}_{4}^{\times}$. The quadrinomial

$$
C(X)=X^{3}+\omega X^{\left(2^{i}+1\right) 2^{k}}+\omega^{2} X^{3 \cdot 2^{m}}+X^{2^{m}\left(2^{i}+1\right) 2^{k}},
$$

is APN on $\mathbb{F}_{2^{n}}$ in the following cases:\\
i) $i=m-2$ and $k$ is even,\\
ii) $i \equiv(m-2)^{-1}(\bmod n)$ and $k$ is even,\\
iii) $i=m$ and $k$ is even,\\
iv) $i=n-1$ and $k$ is even,\\
v) $i=m+2$ and $k$ is odd,\\
vi) $i \equiv(m+2)^{-1}(\bmod n)$ and $k$ is odd,\\
vii) $i=m$ and $k$ is odd.

Now let, as before, $X=x \omega+y$ for $X \in \mathbb{F}_{2^{n}}$, where $x, y \in \mathbb{F}_{2^{m}}$. Then, straightforward computation shows, when $k$ is even,\\
$C_{\text {even }}(x, y)=\omega\left(x^{3}+x^{2} y+y^{3}\right)+\omega^{2}\left(x^{2^{i}+1}+x^{2^{i}} y+y^{2^{i}+1}\right)^{2^{k}}$, and when $k$ is odd,\\
$C_{\text {odd }}(x, y)=\omega\left(x^{3}+x^{2} y+y^{3}\right)+\omega^{2}\left(x^{2^{i}+1}+x y^{2^{i}}+y^{2^{i}+1}\right)^{2^{k}}$.\\
Interchanging ( $x, y$ ) and dividing by $\omega$ lead to EA-equivalent functions. Thus,

$$
\begin{aligned}
C_{\text {even }}(x, y) & \sim_{\text {EA }}\left[x^{3}+x y^{2}+y^{3},\left(x^{2^{i}+1}+x y^{2^{i}}+y^{2^{i}+1}\right)^{2^{k}}\right] \\
& \sim_{\text {EA }}\left[x^{3}+x y^{2}+y^{3}, x^{2^{i}+1}+x y^{2^{i}}+y^{2^{i}+1}\right]
\end{aligned}
$$

since $L(x, y):(x, y) \mapsto\left(x, y^{2^{k}}\right)$ is an $\mathbb{F}_{2}$-linear permutation. Therefore

$$
\begin{aligned}
C_{\text {even }}(x, y) & \sim_{\text {EA }}\left[(1,0,1,1)_{2},(1,0,1,1)_{2^{i}}\right] \\
& \sim_{\text {EA }}\left[(1,0,1,1)_{2},(1,1,0,1)_{2^{m-i}}\right] .
\end{aligned}
$$

Similarly,

$$
C_{\text {odd }}(x, y) \sim_{\text {EA }}\left[(1,0,1,1)_{2},(1,1,0,1)_{2^{i}}\right] .
$$

Now we can observe (see Section III for the family $\mathcal{C}$ ):\\
Proposition VI.2: The APN function $C(X)$ in Theorem VI. 1 is equivalent to:\\
i) $C \sim_{\text {ea }}\left[(1,0,1,1)_{2},(1,1,0,1)_{2^{2}}\right] \in \mathcal{F}_{1}$,\\
ii) $C \sim_{\text {EA }}\left[(1,0,1,1)_{2^{(m+1) / 2}},(1,1,0,1)_{2}\right] \in \mathcal{F}_{1}$\\
iii) $C \sim_{\mathrm{EA}}\left[x^{3}+x y^{2}+y^{3}, x y\right] \in \mathcal{C}$,\\
iv) $C \sim_{\mathrm{EA}} X^{3}$,\\
v) $C \sim_{\text {eA }}\left[(1,0,1,1)_{2},(1,1,0,1)_{2^{2}}\right] \in \mathcal{F}_{1}$,\\
vi) $C \sim_{\text {EA }}\left[(1,0,1,1)_{2^{(m+1) / 2}},(1,1,0,1)_{2}\right] \in \mathcal{F}_{1}$\\
vii) $C \sim_{\text {ea }}\left[x^{3}+x y^{2}+y^{3}, x y\right] \in \mathcal{C}$.

Proof: i) Clear.\\
ii) Note that, if $i \equiv(m-2)^{-1}(\bmod n)$, then $i \equiv-1 / 2 (\bmod m)$. Interchanging $(x, y)$ gives the result.\\
iii) When $i=m, x^{2^{i}+1}+x y^{2^{i}}+y^{2^{i}+1} \sim_{\text {EA }} x y$.\\
iv) The $\mathbb{F}_{2}$-linear map $(x, y) \mapsto(x, x+y)$ is a permutation. Thus $\left[(1,0,1,1)_{2},(1,1,0,1)_{2}\right] \sim_{\text {EA }} \left[(1,0,1,1)_{2},(0,1,1,0)_{2}\right] \quad \sim_{\text {ea }} \quad X^{3}$ as shown in the Introduction.\\
v) Clear.\\
vi) Note that, if $i \equiv(m+2)^{-1}(\bmod n)$, then $i \equiv 1 / 2 (\bmod m)$. Interchanging $(x, y)$ gives the result.\\
vii) Identical to the case iii).

The authors of [7] observed that their construction gives new functions in the cases i), ii), v), and vi). They also observed that it reduces to known families in the other cases iii), iv), and vii), when $m=5$. We have given a general proof here. Thus, the set of new functions in family $\mathcal{B} \mathcal{H} \mathcal{K}$ is now contained in family $\mathcal{F}_{1}$. Those functions in family $\mathcal{F}_{1}$ arise precisely when $i \in\{ \pm 1, \pm 1 / 2\}$ modulo $m$, and $m$ is odd. The only thing we now need to show is that $\mathcal{F}_{1}$ contains new functions which are not in $\mathcal{B} \mathcal{H} \mathcal{K}$. Note that the family $\mathcal{F}_{1}$ does not require $m$ to be odd. This is enough to show that family $\mathcal{F}_{1}$ is new. Nevertheless, we will show that it contains new functions also when $m$ is odd. When $m=5, \mathcal{F}_{1}$ does not contain a new function as $\{1,2,3,4\} \equiv\{ \pm 1, \pm 1 / 2\}(\bmod 5)$, see Table I. When $m=7, \mathcal{F}_{1}$ does contain a new function, see Table II (the symbol † denotes functions which are in $\mathcal{B} \mathcal{H} \mathcal{K}$ ). Tables I and II list the values of the EA-invariant $\mathrm{N}_{F}$ for the

\begin{table}[h]
\begin{center}
\captionsetup{labelformat=empty}
\caption{TABLE I\\
EA-Invariants $N_{F}$ for the Families $\mathcal{F}_{0}, \mathcal{F}_{1}, \mathcal{F}_{2}$ and the Function $\mathcal{F}_{3}$ on $\mathbb{F}_{2^{10}}$}
\begin{tabular}{|l|l|l|}
\hline
Family & $i$ & $\mathrm{~N}_{F}$ \\
\hline
$\mathcal{T}$ & 1,4 & $[341,6820,3720,31]$ \\
 & 2,3 & $[341,6820,3565]$ \\
\hline
$\mathcal{F}_{1}$ & 1,4 & $[341,6820,3565] \dagger$ \\
 & 2,3 & $[341,6820,3720,31] \dagger$ \\
\hline
$\mathcal{F}_{2}$ & 1,4 & $[341,6820,4030,62,1]$ \\
 & 2,3 & $[341,6820,3720,62,1]$ \\
\hline
$\mathcal{F}_{3}$ & n/a & $[341,6324,2573,62,2]$ \\
\hline
\end{tabular}
\end{center}
\end{table}

\begin{table}[h]
\begin{center}
\captionsetup{labelformat=empty}
\caption{TABLE II\\
EA-Invariants $N_{F}$ For Families $\mathcal{F}_{0}, \mathcal{F}_{1}$ and $\mathcal{F}_{2}$ on $\mathbb{F}_{2^{14}}$}
\begin{tabular}{|l|l|l|}
\hline
Family & $i$ & $\mathrm{N}_{F}$ \\
\hline
\multirow{3}{*}{} & 1,6 & [5461, 1681988, 13313156, 436626] \\
\hline
 & 2,5 & [5461, 1681988, 13267817, 401828] \\
\hline
 & 3,4 & [5461, 1681988, 13290042, 428625] \\
\hline
\multirow[t]{3}{*}{$\mathcal{F}_{1}$} & 1,6 & [5461, 1681988, 13250164, 394843] † \\
\hline
 & 2,5 & [5461, 1681988, 13238480, 398399] \\
\hline
 & 3,4 & [5461, 1681988, 13286867, 438531] † \\
\hline
\multirow[t]{3}{*}{$\mathcal{F}_{2}$} & 1,6 & [5461, 1681988, 13219303, 413004, 2667, 127, 1] \\
\hline
 & 2,5 & [5461, 1681988, 13293725, 430784, 2667, 127, 1] \\
\hline
 & 3,4 & [5461, 1681988, 13290423, 418084, 2667, 127, 1] \\
\hline
\end{tabular}
\end{center}
\end{table}

\begin{table}[h]
\begin{center}
\captionsetup{labelformat=empty}
\caption{TABLE III\\
EA-Invariants $N_{F}$ For Known Quadratic APN Functions on $\mathbb{F}_{2}{ }^{10}$}
\begin{tabular}{|l|l|}
\hline
Function $F$ & $\mathrm{N}_{F}$ \\
\hline
$x^{3}$ & [341, 6820, 5115, 341, 11] \\
\hline
$x^{9}$ & [341, 6820, 5115, 341, 11] \\
\hline
$x^{6}+x^{33}+u^{31} x^{192}$ & [341, 6820, 3720, 31] \\
\hline
$x^{33}+x^{72}+u^{31} x^{258}$ & [341, 6820, 3720, 31] \\
\hline
$x^{3}+\operatorname{tr}\left(x^{9}\right)$ & [341, 6820, 4215, 66, 1] \\
\hline
$x^{3}+u^{341} x^{36}$ & [341, 6820, 4400] \\
\hline
$x^{3}+u^{1022} \operatorname{tr}\left(u^{3} x^{9}\right)$ & [341, 6820, 4250, 66, 1] \\
\hline
$x^{57}$ & N/A \\
\hline
$x^{339}$ & N/A \\
\hline
\end{tabular}
\end{center}
\end{table}

families $\mathcal{T}, \mathcal{F}_{1}, \mathcal{F}_{2}$ and the function $\mathcal{F}_{3}$ when $m=5$ and $m=7$ respectively.

\section*{B. Family $\mathcal{F}_{2}$ and the Function $\mathcal{F} 3$}
Table III lists all known APN functions up to CCZequivalence and the known quadratic APN functions (at the time of writing) $F$ on $\mathbb{F}_{2^{10}}$ (as listed in [7] except families $\mathcal{B} \mathcal{H} \mathcal{K}$ and $\mathcal{T}$ ) and their corresponding sequence $\mathrm{N}_{F}$. Due to a result of Dempwolff [21], and Yoshiara [22], a quadratic APN function is CCZ-equivalent to a power function on $\mathbb{F}_{2^{n}}$, if and only if it is EA-equivalent to a Gold APN function $x^{2^{2}+1}$ with $\operatorname{gcd}(i, n)=1$. The only known non-quadratic APN functions on $\mathbb{F}_{2^{10}}$ are CCZ-equivalent to power functions (see Table III). Thus, when $m=5$, it is enough to show that the quadratic functions of families $\mathcal{T}, \mathcal{F}_{1}, \mathcal{F}_{2}$ and the function $\mathcal{F}_{3}$ are not EA-equivalent to the known quadratic APN functions. Now it is clear that any function in family $\mathcal{F}_{2}$ or $\mathcal{F}_{3}$ is not CCZequivalent to a known APN function.

All of the new functions on $\mathbb{F}_{2^{10}}$ and $\mathbb{F}_{2^{14}}$ have the same Walsh spectrum as the Gold functions i.e., $\widehat{F}(U, V) \in \left\{0, \pm 2^{n / 2}, \pm 2^{1+n / 2}\right\}$ and they are not equivalent to a permutation (by a computer search).

\section*{Appendix}
Lemma A.1: We have the following facts.\\
i) If $u \in \mathbb{F}_{4} \backslash \mathbb{F}_{2}$, then the common solutions of $\mathrm{E}_{f}^{u}(x, y)=0=\mathrm{E}_{g}^{u}(x, y)$ are $(x, y) \in \mathbb{F}_{2} \times\{0\}$.\\
ii) The polynomial $\lambda_{u}^{\prime}(y)$ is (see (10))

$$
\begin{aligned}
\lambda_{u}^{\prime}(y) & =u^{q^{2}}\left(\phi_{q}(u)\right)^{2} y^{q^{2}} \\
& +u^{q}\left(\phi_{q^{2}}(u+1)\right)^{2} y^{q} \\
& +u\left(\phi_{q}(u)\right)^{2 q} y
\end{aligned}
$$

iii) The polynomial

$$
\pi_{u}^{\prime}(x)=\phi_{q}(u) x^{q+1}+\phi_{q^{2}}(u+1) x+\left(\phi_{q}(u)\right)^{q}
$$

satisfies

$$
\pi_{u}^{\prime}(x)=\left(\epsilon_{3}^{\prime} x+\epsilon_{4}^{\prime}\right)^{q+1} \phi_{q}\left(\frac{\epsilon_{1}^{\prime} x+\epsilon_{2}^{\prime}}{\epsilon_{3}^{\prime} x+\epsilon_{4}^{\prime}}\right)
$$

with

$$
\left(\begin{array}{ll}
\epsilon_{1}^{\prime} & \epsilon_{2}^{\prime} \\
\epsilon_{3}^{\prime} & \epsilon_{4}^{\prime}
\end{array}\right)=\left(\begin{array}{cc}
1 & (u+1)^{q} \\
(u+1) & u^{q}
\end{array}\right) .
$$

Proof: i) Assume $y \neq 0$. By Lemma IV. 4 (ii) and $\mathrm{E}_{g}^{u}(x, y)=0$ we get $y^{q^{2}-1}=u^{2}$, hence $u^{2}$ must be a cube. On the other hand, $u^{2}$ is a primitive third root of unity, hence is of the form $u^{2}=\alpha^{\frac{2^{m}-1}{3}}$, where $\alpha$ is a generator of $\mathbb{F}_{2^{m}}^{*}$. But now, $u^{2}$ is a cube if and only if $3 \mid\left(2^{m}-1\right) / 3$ if and only if $9 \mid 2^{m}-1$ if and only if $2^{m} \equiv 1(\bmod 9)$ if and only if $6 \mid m$, which is impossible since $\operatorname{gcd}(3, m)=1$.\\
ii) We need to get to (11) from (10).\\
a) The equality for the coefficient of $y^{q^{2}}$ is

$$
\begin{aligned}
u^{q^{2}} \phi_{q}(u) & =(u+1)^{q} \phi_{q^{2}}(u+1) \\
& +\left[\phi_{q}(u)\right]^{q} .
\end{aligned}
$$

We get

$$
\begin{aligned}
& {\left[\phi_{q^{2}}(u+1)\right](u+1)^{q}+\left[\phi_{q}(u)\right]^{q}} \\
& =\left(u^{q^{2}+1}+u^{q^{2}}+1\right)\left(u^{q}+1\right)+u^{q^{2}+q}+u^{q}+1 \\
& =u^{q^{2}+q+1}+u^{q^{2}+1}+u^{q^{2}} \\
& =u^{q^{2}}\left(u^{q+1}+u+1\right)=u^{q^{2}} \phi_{q}(u) .
\end{aligned}
$$

b) The equality for the coefficient of $y^{q}$ is

$$
u^{q} \phi_{q^{2}}(u+1)=(u+1)\left[\phi_{q}(u)\right]^{q}+\phi_{q}(u) .
$$

We get

$$
\begin{aligned}
& (u+1)\left[\phi_{q}(u)\right]^{q}+\phi_{q}(u) \\
& =(u+1)\left(u^{q^{2}+q}+u^{q}+1\right)+u^{q+1}+u+1 \\
& =u^{q^{2}+q+1}+u^{q^{2}+q}+u^{q} \\
& =u^{q}\left(u^{q^{2}+1}+u^{q^{2}}+1\right)=u^{q} \phi_{q^{2}}(u+1) .
\end{aligned}
$$

c) The equality for the coefficient of $y$ is

$$
u\left[\phi_{q}(u)\right]^{q}=(u+1)^{q^{2}} \phi_{q}(u)+\phi_{q^{2}}(u+1) .
$$

We get

$$
\begin{aligned}
& (u+1)^{q^{2}} \phi_{q}(u)+\phi_{q^{2}}(u+1) \\
& =\left(u^{q^{2}}+1\right)\left(u^{q+1}+u+1\right)+u^{q^{2}+1}+u^{q^{2}}+1 \\
& =u^{q^{2}+q+1}+u^{q+1}+u \\
& =u\left(u^{q+1}+u+1\right)^{q}=u\left[\phi_{q}(u)\right]^{q} .
\end{aligned}
$$

iii) We have,

$$
\begin{aligned}
& \left(\epsilon_{3}^{\prime} x+\epsilon_{4}^{\prime}\right)^{q+1} \phi_{q}\left(\frac{\epsilon_{1}^{\prime} x+\epsilon_{2}^{\prime}}{\epsilon_{3}^{\prime} x+\epsilon_{4}^{\prime}}\right)=\left(x+(u+1)^{q}\right)^{q+1} \\
& +\left(x+(u+1)^{q}\right)\left((u+1)^{q} x^{q}+u^{q^{2}}\right) \\
& +\left((u+1) x+u^{q}\right)^{q+1} \\
= & \left(x^{q+1}+x^{q}(u+1)^{q}+x(u+1)^{q^{2}}+(u+1)^{q^{2}+q}\right) \\
& +\left(x^{q+1}(u+1)^{q}+x^{q}(u+1)^{2 q}+x u^{q^{2}}+u^{q^{2}}(u+1)^{q}\right) \\
& +\left(x^{q+1}(u+1)^{q+1}+x^{q}(u+1)^{q} u^{q}\right. \\
& \left.+x u^{q^{2}}(u+1)+u^{q^{2}+q}\right) \\
= & x^{q+1}\left(1+(u+1)^{q}+(u+1)^{q+1}\right) \\
& +x^{q}\left((u+1)^{q}\left(1+(u+1)^{q}+u^{q}\right)\right) \\
& +x\left(u^{q^{2}+1}+u^{q^{2}}+1\right)+u^{q^{2}+q}+u^{q^{2}} \\
& +u^{q}+1+u^{q^{2}+q}+u^{q^{2}}+u^{q^{2}+q} \\
= & \phi_{q}(u) x^{q+1}+\phi_{q^{2}}(u+1) x+\left(\phi_{q}(u)\right)^{q}=\pi_{u}^{\prime}(x) .
\end{aligned}
$$

Lemma A.2: i) The polynomial $\tau_{u}^{\prime}(y)$ is (see (14))

$$
\begin{aligned}
\tau_{u}^{\prime}(y) & =\frac{u^{q^{3}} \phi_{q^{2}}(u+1)}{\left(\phi_{q}(u)\right)^{q^{2}} \psi_{q^{3}}(u)} y^{q^{3}}+\frac{u^{q^{2}}\left(\phi_{q^{2}}(u+1)\right)^{q}}{\left(\phi_{q}(u)\right)^{q^{2}+q}} y^{q^{2}} \\
& +\frac{u^{q} \phi_{q^{2}}(u+1)}{\left(\phi_{q}(u)\right)^{q+1}} y^{q}+\frac{u\left(\phi_{q^{2}}(u+1)\right)^{q}}{\phi_{q}(u) \psi_{q^{3}}(u)} y
\end{aligned}
$$

ii) We have

$$
\tau_{u}(y)=\lambda_{u}(y)+\frac{\left(\lambda_{u}(y)\right)^{q}}{\left(\phi_{q^{2}}(u+1)\right)^{q-1}\left(\phi_{q}(u)\right)^{q^{2}-1}} .
$$

iii) For all $u \in \mathbb{F}_{2^{m}} \backslash \mathbb{F}_{4}$, we have

$$
\lambda_{u}\left(u^{2}+u+1\right)=\phi_{q^{2}}(u+1)\left(\phi_{q}(u)\right)^{q+1},
$$

where

$$
\lambda_{u}(y)=\left(\phi_{q}(u)\right)^{2} y^{q^{2}}+\left(\phi_{q^{2}}(u+1)\right)^{2} y^{q}+\left(\phi_{q}(u)\right)^{2 q} y .
$$

iv) For $u \in \mathbb{F}_{2^{m}}$, we have

$$
\mu_{u}\left(\frac{u^{2}+u+1}{u}\right)=\frac{1}{u^{q}}+\frac{1}{u}+1,
$$

where

$$
\mu_{u}(y)=\frac{(u+1) y^{q}+y}{\phi_{q}(u)} .
$$

Proof: i) From (14) we get to $\tau_{u}^{\prime}(y)$ as follows:\\
a) The coefficient of $y^{q^{3}}$ is verified by

$$
\begin{aligned}
& \left(u^{q^{3}}+u\right)(u+1)^{q^{2}}+u\left(\phi_{q}(u)\right)^{q^{2}} \\
= & u^{q^{2}(q+1)}+u^{q^{3}}+u^{q^{2}+1}+u+u^{q^{2}(q+1)+1} \\
& +u^{q^{2}+1}+u \\
= & u^{q^{3}}\left(u^{q^{2}+1}+u^{q^{2}}+1\right) \\
= & u^{q^{3}} \phi_{q^{2}}(u+1) .
\end{aligned}
$$

b) For the verification of the coefficient of $y^{q^{2}}$, use the equality in Lemma A. 1 (ii) (b) (taken to $q$ th power).\\
c) For the verification of the coefficient of $y^{q}$, use the equality in Lemma A. 1 (ii) (b).\\
d) The coefficient of $y$ is verified as follows.

$$
u^{q^{3}} \phi_{q}(u)+u^{q^{3}}+u=u\left(\phi_{q^{2}}(u+1)\right)^{q}
$$

if and only if

$$
u^{q^{3}}\left(\phi_{q}(u)+1\right)=u\left(\left(\phi_{q^{2}}(u+1)\right)^{q}+1\right)
$$

if and only if

$$
u^{q^{3}+q+1}+u^{q^{3}+1}=u^{q^{3}+q+1}+u^{q^{3}+1}
$$

ii) The coefficients of the terms $y$ and $y^{q^{3}}$ are clear. Consider

$$
\begin{aligned}
& t(y)=\frac{\tau(y)}{y}, \text { and } \\
& l(y)=\frac{\lambda_{u}(y)}{y}+\frac{\left(\lambda_{u}(y)\right)^{q}}{C y}
\end{aligned}
$$

Now we will show that the $y^{q-1}$ coefficients of $l\left(y \phi_{q^{2}}(u+1)\right)$ and $t\left(y \phi_{q^{2}}(u+1)\right)$ are the same. Indeed,

$$
\begin{aligned}
& \left(\phi_{q^{2}}(u+1)\right)^{2}\left(\phi_{q^{2}}(u+1)\right)^{q-1} \\
& +\frac{\left(\phi_{q}(u)\right)^{2 q^{2}}\left(\phi_{q^{2}}(u+1)\right)^{q-1}}{\left(\phi_{q^{2}}(u+1)\right)^{q-1}\left(\phi_{q}(u)\right)^{q^{2}-1}} \\
& =\left(\phi_{q^{2}}(u+1)\right)^{q+1}+\left(\phi_{q}(u)\right)^{q^{2}+1} \\
& =u^{q^{3}+q^{2}+q}+u^{q^{3}+q}+u^{q^{3}}+u^{q^{2}+q+1}+u^{q+1}+u \\
& =\left(u^{q^{2}+q}+u^{q}+1\right)\left(u^{q^{3}}+u\right) \\
& =\left(\phi_{q}(u)\right)^{q} \psi_{q^{3}}(u)
\end{aligned}
$$

Similarly, the $y^{q^{2}-1}$ coefficients of $l\left(y \phi_{q}(u)\right)$ and $t\left(y \phi_{q}(u)\right)$ are the same. That is,

$$
\begin{aligned}
& \left(\phi_{q}(u)\right)^{2}\left(\phi_{q}(u)\right)^{q^{2}-1}+\frac{\left(\phi_{q^{2}}(u+1)\right)^{2 q}\left(\phi_{q}(u)\right)^{q^{2}-1}}{\left(\phi_{q^{2}}(u+1)\right)^{q-1}\left(\phi_{q}(u)\right)^{q^{2}-1}} \\
& =\left(\phi_{q}(u)\right)^{q^{2}+1}+\left(\phi_{q^{2}}(u+1)\right)^{q+1}
\end{aligned}
$$

which is identical to the $y^{q-1}$ case.\\
iii) The polynomial equality

$$
\begin{aligned}
\phi_{q^{2}}(u+1)\left(\phi_{q}(u)\right)^{q+1} & =\left(\phi_{q}(u)\right)^{2}\left(u^{2}+u+1\right)^{q^{2}} \\
& +\left(\phi_{q^{2}}(u+1)\right)^{2}\left(u^{2}+u+1\right)^{q} \\
& +\left(\phi_{q}(u)\right)^{2 q}\left(u^{2}+u+1\right)
\end{aligned}
$$

can be written as

$$
\begin{aligned}
\sum_{i \in I_{1}} u^{i} \sum_{i \in I_{2}} u^{i} \sum_{i \in I_{3}} u^{i} & =\sum_{i \in J_{1}} u^{i} \sum_{i \in J_{2}} u^{i} \\
& +\sum_{i \in J_{3}} u^{i} \sum_{i \in J_{4}} u^{i} \\
& +\sum_{i \in J_{5}} u^{i} \sum_{i \in J_{6}} u^{i}
\end{aligned}
$$

where

$$
\begin{aligned}
I_{1} & =\left[q^{2}+1, q^{2}, 0\right], \\
I_{2} & =\left[q^{2}+q, q, 0\right], \\
I_{3} & =[q+1,1,0], \\
J_{1} & =[2 q+2,2,0], \\
J_{2} & =\left[2 q^{2}, q^{2}, 0\right], \\
J_{3} & =\left[2 q^{2}+2,2 q^{2}, 0\right], \\
J_{4} & =[2 q, q, 0], \\
J_{5} & =\left[2 q^{2}+2 q, 2 q, 0\right], \\
J_{6} & =[2,1,0] .
\end{aligned}
$$

Consider the direct sums of these multisets:

$$
\begin{aligned}
M_{1} & =\left[q^{2}+1, q^{2}, 0\right] \\
& +\left[q^{2}+q, q, 0\right] \\
& +[q+1,1,0] \\
M_{2} & =[2 q+2,2,0]+\left[2 q^{2}, q^{2}, 0\right] \\
& \cup\left[2 q^{2}+2,2 q^{2}, 0\right]+[2 q, q, 0] \\
& \cup\left[2 q^{2}+2 q, 2 q, 0\right]+[2,1,0] .
\end{aligned}
$$

which can be written in a more readable vector-form

$$
\begin{aligned}
M_{1} & =[101,100,000] \\
& +[110,010,000] \\
& +[011,001,000] \\
M_{2} & =([022,002,000]+[200,100,000]) \\
& \cup([202,200,000]+[020,010,000]) \\
& \cup([220,020,000]+[002,001,000]),
\end{aligned}
$$

where each element $a_{0} a_{1} a_{2}$ is viewed as a vector $\left(a_{0}, a_{1}, a_{2}\right)$ in $\mathbb{Z}^{3}$, and correspond to $a_{2} q^{2}+a_{1} q+a_{0}$, e.g., $101=(1,0,1) \in \mathbb{Z}^{3}$ and corresponds to $q^{2}+1$. We need to show that the set of those elements in $M_{1}$ and $M_{2}$ that have odd repetition numbers (say $N_{1}$ and $N_{2}$ ) are equal. The fact that $N_{1}=N_{2}$ can be verified by a simple Magma script. We have

$$
\begin{aligned}
N_{1}=N_{2}=\{ & 001,000,222,221,122, \\
& 102,210,100,212,021,010\} .
\end{aligned}
$$

iv) The claim is simply verified by checking

$$
\frac{(u+1) \frac{\left(u^{2}+u+1\right)^{q}}{u^{q}}+\frac{u^{2}+u+1}{u}}{\phi_{q}(u)}=\frac{1}{u^{q}}+\frac{1}{u}+1 .
$$

This is true since we have

$$
\begin{aligned}
& \frac{u(u+1)\left(u^{2}+u+1\right)^{q}+u^{q}\left(u^{2}+u+1\right)}{u^{q+1}\left(u^{q+1}+u+1\right)}= \\
& \frac{u^{q+1}+u^{q}+u}{u^{q+1}}
\end{aligned}
$$

if and only if

$$
\begin{aligned}
& \left(u^{2}+u\right)\left(u^{2 q}+u^{q}+1\right)+u^{q}\left(u^{2}+u+1\right)= \\
& \left(u^{q+1}+u+1\right)\left(u^{q+1}+u^{q}+u\right)
\end{aligned}
$$

if and only if

$$
\begin{aligned}
& u^{2 q+2}+u^{2}+u^{2 q+1}+u+u^{q}= \\
& u^{2 q+2}+u^{2 q+1}+u^{2}+u^{q}+u .
\end{aligned}
$$

\section*{Acknowledgment}
The authors thank an editor and a reviewers for extensive reviews which improved the presentation of the paper.

\section*{References}
[1] K. Browning, J. Dillon, M. McQuistan, and A. Wolfe, "An APN permutation in dimension six," in Finite Fields: Theory and Applications-FQ9 (Contemporary Mathematics), vol. 518. Providence, RI, USA: AMS, 2010, pp. 33-42.\\[0pt]
[2] C. Carlet, P. Charpin, and V. Zinoviev, "Codes, bent functions and permutations suitable for DES-like cryptosystems," Des., Codes Cryptogr., vol. 15, no. 2, pp. 125-156, Nov. 1998.\\[0pt]
[3] S. Yoshiara, "Equivalences of quadratic APN functions," J. Algebr. Combinatorics, vol. 35, no. 3, pp. 461-475, May 2012.\\[0pt]
[4] F. Göloğlu and J. Pavlu̇, "On CCZ-inequivalence of some families of almost perfect nonlinear functions to permutations," Cryptogr. Commun., vol. 13, no. 3, pp. 377-391, May 2021.\\[0pt]
[5] R. Gold, "Maximal recursive sequences with 3 -valued recursive crosscorrelation functions (Corresp.)," IEEE Trans. Inf. Theory, vol. IT-14, no. 1, pp. 154-156, Jan. 1968.\\[0pt]
[6] X.-D. Hou, "Affinity of permutations of $\mathbb{F}_{2}^{n}$," Discrete Appl. Math., vol. 154, no. 2, pp. 313-325, 2006.\\[0pt]
[7] L. Budaghyan, T. Helleseth, and N. Kaleyski, "A new family of APN quadrinomials," IEEE Trans. Inf. Theory, vol. 66, no. 11, pp. 7081-7087, Nov. 2020.\\[0pt]
[8] T. P. Berger, A. Canteaut, P. Charpin, and Y. Laigle-Chapuy, "On almost perfect nonlinear functions over $m m b F_{n}^{2}$," IEEE Trans. Inf. Theory, vol. 52, no. 9, pp. 4160-4170, Sep. 2006.\\[0pt]
[9] K. Nyberg, "S-boxes and round functions with controllable linearity and differential uniformity," in Proc. Int. Workshop Fast Softw. Encryption in Lecture Notes in Computer Science, vol. 1008, B. Preneel, Ed. Cham, Switzerland: Springer, 1994, pp. 111-130.\\[0pt]
[10] S. Abhyankar, "Projective polynomials," Proc. Amer. Math. Soc., vol. 125, no. 6, pp. 1643-1650, 1997.\\[0pt]
[11] A. W. Bluher, "On $x^{q+1}+a x+b$," Finite Fields Appl., vol. 10, no. 3, pp. 285-305, 2004.\\[0pt]
[12] C. Carlet, Boolean Functions for Cryptography Coding Theory. Cambridge, U.K.: Cambridge Univ. Press, 2021.\\[0pt]
[13] R. Lidl and H. Niederreiter, "Finite fields and their applications," in Handbook Algebra, vol. 1. Amsterdam, The Netherlands: North-Holland, 1996.\\[0pt]
[14] Y. Zhou and A. Pott, "A new family of semifields with 2 parameters," Adv. Math., vol. 234, pp. 43-60, Feb. 2013.\\[0pt]
[15] H. Taniguchi, "On some quadratic APN functions," Des., Codes Cryptogr., vol. 87, no. 9, pp. 1973-1983, Sep. 2019.\\[0pt]
[16] C. Carlet, "Relating three nonlinearity parameters of vectorial functions and building APN functions from bent functions," Des., Codes Cryptogr., vol. 59, nos. 1-3, pp. 89-109, 2011.\\[0pt]
[17] L. Budaghyan, C. Carlet, and G. Leander, "Two classes of quadratic APN binomials inequivalent to power functions," IEEE Trans. Inf. Theory, vol. 54, no. 9, pp. 4218-4229, Sep. 2008.\\[0pt]
[18] L. Budaghyan, C. Carlet, and G. Leander, "Constructing new APN functions from known ones," Finite Fields Appl., vol. 15, pp. 150-159, Apr. 2009.\\[0pt]
[19] L. Carlitz, "Explicit evaluation of certain exponential sums," Math. Scand., vol. 44, no. 1, pp. 5-16, 1979.\\[0pt]
[20] P. Sziklai, "On subsets of $\mathrm{GF}\left(\mathrm{q}^{2}\right)$ with dth power differences," Discrete Math., vols. 208-209, pp. 547-555, Oct. 1999.\\[0pt]
[21] U. Dempwolff, "CCZ equivalence of power functions," Des., Codes Cryptogr., vol. 86, no. 3, pp. 665-692, Mar. 2018.\\[0pt]
[22] S. Yoshiara, "Equivalences of power APN functions with power or quadratic APN functions," J. Algebr. Combinatorics, vol. 44, no. 3, pp. 561-585, Nov. 2016.


\end{document}